\documentclass[11pt]{article}
\usepackage[T1]{fontenc}
\usepackage[utf8]{inputenc}
\usepackage{lmodern,microtype,amsmath,amssymb,amsthm,mathtools}
\usepackage[margin=1in]{geometry}
\usepackage{booktabs,array,longtable,enumitem}
\usepackage[round,authoryear]{natbib}
\usepackage[colorlinks=true,linkcolor=blue!50!black,citecolor=blue!50!black,urlcolor=blue!50!black]{hyperref}
\usepackage{xcolor}
\hypersetup{pdftitle={The cost of coalition withholding: resource dimension and coalition size},pdfauthor={Maxwell Flanagan},pdfsubject={Specialist correctness and priority-review draft}}
\newcommand{\R}{\mathbb R}
\newcommand{\E}{\mathbb E}
\newcommand{\ind}{\mathbf1}
\newtheorem{theorem}{Theorem}[section]
\newtheorem{lemma}[theorem]{Lemma}
\newtheorem{proposition}[theorem]{Proposition}
\newtheorem{corollary}[theorem]{Corollary}
\theoremstyle{definition}\newtheorem{definition}[theorem]{Definition}
\theoremstyle{remark}\newtheorem{remark}[theorem]{Remark}
\setlength{\emergencystretch}{2em}
\title{The cost of coalition withholding:\\resource dimension and coalition size}
\author{Maxwell Flanagan\\[2pt]\small New York City, United States\\[2pt]
\small Corresponding author: \texttt{maxwell.flanagan@gmail.com}}
\date{24 October 2026\\\small Specialist correctness and priority-review draft; not submitted}
\begin{document}
\maketitle
\begin{abstract}
Owners contribute resource bundles to a common production technology. Payments
must cover every coalition's stand-alone production and resist coordinated
resource withholding, with side payments allowed inside the deviating coalition.
We distinguish payment-only protection from protection that also compensates
the coalition's outside production. For finite-valued monotone concave
homogeneous technologies on $d$ resources, the sharp worst-instance total-funding
factors are $(d+1)/2$ and $1+c(d-1)$ respectively, where
$c(2-e^{-1/c})=1$ and $c\simeq0.5430633898$.
For packing production of activity arity at most $r$, protecting coalitions of
at most $s$ owners admits a full-outside funding factor $rs^{r-1}$,
independently of the owner and resource counts.
A projective construction gives an unrestricted payment-only lower
$K_3(s)\ge\sqrt{s/6}$, with a sharper explicit certificate asymptotic to
$\sqrt{2\pi s}/(3\sqrt e)$. It extends to
$K_r(s)=\Omega_r(s^{(r-2)/(r-1)})$ for every fixed $r\ge3$.
For unit-profit four-vertex paths, complete path blow-ups and their disjoint
unions, an explicit coupled flow selector gives full outside protection with
total-funding factor $1+1/\sqrt2$, independently of owner and resource counts.
All rules and deviations retain fixed actual owners and their full real report
boxes. We give complete proofs, distinguish mixed and pure ownership, and
supply exact finite certificate checks. The arity-three gap to $3s^2$ remains
open. The quantitative claims are submitted for specialist correctness and
priority assessment; the underlying monotonicity axioms, semivalues and
hypergraph ingredients have established antecedents.
\end{abstract}
\smallskip
\noindent\textbf{Keywords:} coalition incentives; resource withholding;
production games; subsidy; sparse activities; cost sharing.
\clearpage
\tableofcontents
\newpage

\section{Introduction and relation to prior work}\label{sec:intro}
An efficient division can protect every coalition at one production profile
and still fail when several owners withhold resources together. This paper
studies the extra funding needed to choose payments consistently across all
feasible reports. Owners may supply several inputs. Their identities,
eligibility and endowment boxes remain fixed, and a protected coalition may
include members whose reports do not change. Transfers inside the coalition
make its total payment the relevant incentive quantity.

Two parameters organize the results. Resource dimension $d$ measures the
number of distinct inputs when every coalition is protected. Activity arity
$r$ instead bounds the inputs used by one production activity, while both the
number of resources and owners may grow; $s$ bounds only the size of coalitions
receiving incentive protection. Static coverage remains unrestricted.
A sponsor's factor $\kappa$ bounds total payments by $\kappa$ times realized
production. The external supplement is therefore $(\kappa-1)$ times
production, not $\kappa$ times production.

The paper joins the sharp dimension frontiers with a new lower-bound
construction for bounded-coalition protection. The dimension results include
mixed-owner upper formulas and a reduction from arbitrary finite-owner rules
for the matching outside-production lower. The sparse upper is a binomial
marginal construction using bounded-rank positive supports. The arity-three
lower uses a constant-size productive seed, lawful symmetry of a fixed
institution, and a ladder of one-owner restorations accompanied by unchanged
protected members. The construction proves a funding obstruction for arbitrary
payment functions, rather than only for a prescribed family of mechanisms.
Appendix~\ref{sec:arity} gives the higher-arity and pure-owner versions in the
same paper. Section~\ref{graph:section} adds a constructive counterpart for a restricted
class of unit-profit resource graphs. Its protected-leaf maximum-flow selector
satisfies all-coalition outside protection with factor $1+1/\sqrt2$ on complete
four-vertex-path blow-ups and disjoint unions, with arbitrary mixed ownership.
This is a scoped graph construction, not an upper bound for arbitrary graphs
or a claim of optimality on that class. The quantitative comparisons and proof
dependencies are summarized in the separate specialist brief.

\paragraph{Demand monotonicity and absence-proofness.}
The aggregate-sum monotonicity used here has a close antecedent in strong group
monotonicity of \citet{MoulinSprumont2004}. Their scalar integer-demand
formulation permits passive members of the protected group. We do not claim
that inequality as a new axiom. Their Proposition~2(ii), in the inspected 2004
author version, also assumes additivity, separability and exact balance; its
impossibility does not directly apply to our unrestricted overfunded rules.
The production class matters as well. \citet{Dogan2013} studies
absence-proofness with coalition outside worth in efficient subgame allocation
schemes. That is a close conceptual comparison to full outside protection.
Its inspected working-paper contract concerns binary absence and efficient
subgames. Whether a published or related result implies a present quantitative
bound is an explicit priority-review question.

\paragraph{Cross-monotonic sharing.}
\citet{ImmorlicaEtAl2008} give quantitative limitations of cross-monotonic
cost and profit sharing. Their profit formulation has compatible overfunding
and core signs. Its cross-monotonicity constrains each incumbent separately
under participation changes, whereas our withholding condition constrains the
sum for a coalition that can include a restoring owner. Their maximum-flow
lower uses the individual constraints. It cannot be transferred solely from
the theorem headline; an additional reduction would be needed. We make no
claim that all related lower-bound techniques have been exhausted.

\paragraph{Population-monotonic matching allocations.}
\citet{XiaoFang2021} study efficient allocation schemes for induced-subgraph
matching games whose players are vertices. Their population monotonicity
requires each incumbent's share to increase when the population expands.
The P4 property in their Lemma~3 rules out the unit-weight path under that
contract. Our graph theorem instead permits subsidized payments, divisible
resource reports and mixed ownership, and protects a deviating coalition's
aggregate payment plus its outside production. It does not impose individual
incumbent population monotonicity. Consequently that obstruction neither
contradicts nor establishes priority for the construction below.

\paragraph{Static subsidies and restricted coalitions.}
The cost of stability of \citet{BachrachEtAl2009} concerns subsidies making a
single cooperative game stable, and already includes projective geometry and
square-root phenomena. Neither ingredient alone is new here. Every packing
profile in this paper already has an efficient static cover by dual prices;
the obstruction is consistency across reports.
In \citet{MeirEtAl2011}, a coalition-size restriction removes static blocking
constraints. Here all coalition coverage constraints are retained, and the
size restriction applies only to incentive comparisons. The parameters cannot
be identified without changing the problem.

\paragraph{Positive supports and semivalues.}
The hierarchy of positive hypergraph representations developed by
\citet{FeigeEtAl2014}, and the bounded-rank loss inequalities in
\citet{EzraEtAl2018}, are antecedents for the sparse upper.
Binomial marginal values and their relation to collusion are also established;
see, for example, \citet{Casajus2011}. We include a direct proof of the exact
upper used here, including full outside protection and mixed ownership, rather
than attributing novelty to the probabilistic value itself. The central-binomial
identity in the lower-bound ladder is a separate arithmetic ingredient.

\paragraph{Review status and scope.}
The proposed contribution is the collection of quantitative funding bounds
under the stated contracts, especially the arbitrary-rule projective ladder and the coupled graph selector.
This draft is not a claim of verified historical priority. The review packet
identifies original versions inspected, missing published-source comparisons,
and additional cost-sharing leads. Internal AI-assisted audits and exact
finite certificates do not replace specialist review of either correctness or
novelty. No empirical pooling, power-reliability or opportunity-cost conclusion
is included.

\section{Model and funding frontiers}\label{sec:model}
An institution fixes a finite set $N$ of actual owners and $d$ resource types.
Owner $i$ has an endowment $\bar b_i\in\R_+^d$ and may report any
$b_i\in[0,\bar b_i]$ coordinatewise. Zero endowment coordinates encode
ineligibility. These owners, coordinates, boxes and the public technology
remain fixed throughout every comparison. Write $b_C=\sum_{i\in C}b_i$,
$B=b_N$, and $\mathcal B=\prod_{i\in N}[0,\bar b_i]$.
Different institutions may have different, arbitrarily large finite owner and
resource counts.

The technology $F:\R_+^d\to\R_+$ is finite-valued, monotone, concave and
homogeneous of degree one, with $F(0)=0$. These properties imply
superadditivity:
\begin{equation}\label{eq:superadditive}
 F(x+y)=2F((x+y)/2)\ge F(x)+F(y).
\end{equation}
No closed-orthant continuity assumption is imposed on the general dimension
class. In the sparse-production sections we specialize to a finite packing
program
\begin{equation}\label{eq:packing}
 F(c)=\max\{\pi\cdot x: Ax\le c,\ x\ge0\},\qquad A\ge0,\quad\pi\ge0.
\end{equation}
Every positive-profit column consumes a positive amount of at least one
resource. Zero-profit columns may be discarded. The optimum is finite and
attained at finite capacities. Its \emph{activity arity} is the maximum number
of positive entries in a positive-profit column. Arity three permits many
resource rows; it does not mean three resources in total.

A rule is an arbitrary map $y:\mathcal B\to\R_+^N$.
It need not be continuous, anonymous, marginal, or linear. All requirements
below concern the same rule at every real report profile. For
$y_C(b)=\sum_{i\in C}y_i(b)$, \emph{coalition coverage} is
\begin{equation}\label{eq:coverage}
 y_C(b)\ge F(b_C)\qquad(C\subseteq N,\ b\in\mathcal B).
\end{equation}
Coverage is required for \emph{every} coalition, without a size restriction.
A total-funding factor $\kappa$ means
\begin{equation}\label{eq:funding}
 \sum_{i\in N}y_i(b)\le\kappa F(B)\qquad(b\in\mathcal B).
\end{equation}
Thus the external subsidy is at most $(\kappa-1)F(B)$; $\kappa$ itself is a
factor on total payments. When $F(B)=0$, nonnegativity and
\eqref{eq:funding} force every payment to vanish. At a positive-production
profile, an owner reporting zero is allowed to receive a positive payment.
There is no deletion-marginal floor or separate dummy-zero axiom.

\begin{definition}[Coalition withholding contracts]\label{def:contracts}
Fix an integer $s\ge1$. For every $C\subseteq N$ with $|C|\le s$, let
$\delta_i\in[0,b_i]$ for $i\in C$ and $\delta_i=0$ outside $C$.
\emph{Transferable coalition withholding monotonicity} (TCWM) requires
\begin{equation}\label{eq:tcwm}
 y_C(b)-y_C(b-\delta)\ge0.
\end{equation}
\emph{Full outside-production protection} (full OUT) requires
\begin{equation}\label{eq:out}
 y_C(b)-y_C(b-\delta)\ge F(\delta_C).
\end{equation}
Protected coalitions may contain unchanged members: $\delta_i=0$ is allowed
for any member of $C$. For unrestricted coalition protection, these
inequalities are required for every $C\subseteq N$.
\end{definition}

Both contracts permit side payments within the protected coalition. TCWM
uses payment-only utility; full OUT additionally allows the coalition to
produce with its pooled withheld resources. Neither is a model of arbitrary
misreports, private opportunity costs, or an endogenous change in the owner
roster. Full OUT implies TCWM. A lower bound for TCWM therefore also applies
to full OUT, whereas a TCWM construction need not satisfy full OUT.

For an institution $I$, let $k_s(I)$ and $k_s^{\mathrm{OUT}}(I)$ be the
infima of its feasible factors under these respective contracts and
\eqref{eq:coverage}. Define
\begin{equation}\label{eq:frontiers}
 K_r(s)=\sup_{I:\,\text{arity}(I)\le r} k_s(I),\qquad
 K_r^{\mathrm{OUT}}(s)=\sup_{I:\,\text{arity}(I)\le r}
 k_s^{\mathrm{OUT}}(I).
\end{equation}
The supremum ranges over all finite packing institutions, including their
actual eligibility and endowment boxes. A uniform upper constructs one rule
for each entire institution; a lower constructs an institution on which
\emph{every} rule has an expensive feasible report. The infimum convention
does not assume optimal rules are attained. We use $D_d$ and
$D_d^{\mathrm{OUT}}$ for the corresponding worst-instance factors for general
$d$-resource technologies when \emph{every} coalition receives protection.

\begin{remark}[Static feasibility versus cross-report incentives]\label{rem:static}
For a packing technology, at each fixed profile an optimal nonnegative dual
price $w$ gives payments $w\cdot b_i$. Dual feasibility implies
$w\cdot b_C\ge F(b_C)$, and their sum is $F(B)$.
Thus each individual report profile admits an efficient coalition cover.
The funding problem concerns selecting payments consistently across profiles.
Neither the static core nor an integrality-gap calculation by itself proves
our cross-report lower bounds.
\end{remark}


\section{The universal resource-dimension frontier}
\label{dim:section}

Here the technology may be any finite-valued function
$F:\R_+^d\to\R_+$ that is coordinatewise nondecreasing, concave, and
positively homogeneous of degree one. In particular, $F(0)=0$.
Continuity between different support faces is not assumed.
The owners, their eligibility, and their finite coordinatewise endowment
bounds are fixed throughout each mechanism. All real reports in those
boxes, including zero reports, are in its domain.

Write $D_d$ for the supremum, over these technologies and fixed finite
owner models, of the infimal funding factor for nonnegative,
coalition-covering TCWM payments. Define $D_d^{\mathrm{OUT}}$ in the same
way with full coalition outside-production protection. Write $Y_S(b)=\sum_{i\in S}y_i(b)$. Thus coverage means
$Y_S(b)\ge F(\sum_{i\in S}b_i)$ for every coalition, and full OUT means
\begin{equation}\label{dim:out-contract}
 Y_S(b)-Y_S(b-w_S)\ge
 F\left(\sum_{i\in S}w_i\right),
 \qquad 0\le w_i\le b_i\quad(i\in S).
\end{equation}
TCWM replaces the right-hand side by zero. A coalition may contain members
whose reports do not change. A funding factor $K$ means
$\sum_i y_i(b)\le K F(\sum_i b_i)$ at every report, including those with
zero production. These definitions do not presume attainment of a
particular model's infimal factor.

\begin{theorem}[Sharp universal dimension factors]\label{dim:main}
For every integer $d\ge1$,
\begin{equation}\label{dim:frontiers}
 D_d=\frac{d+1}{2},\qquad
 D_d^{\mathrm{OUT}}=1+c(d-1),
 \qquad c\bigl(2-e^{-1/c}\bigr)=1,
\end{equation}
where $c$ is the unique solution in $(1/2,1)$
($c\simeq0.5430633898$).
The upper bounds use explicit payment formulas for arbitrary fixed mixed
owners. At a report with exactly $k\ge1$ active resource coordinates they
give the stronger respective factors $(k+1)/2$ and $1+c(k-1)$; at the zero
report they pay zero. The matching lower bounds already hold for
$F(C)=\min_{1\le j\le d}C_j$ on fixed finite endowment boxes, uniformly
over the number of actual owners.
\end{theorem}

The sharpness is a worst-case statement across finite owner models.
It does not assert the same lower bound for every technology or every fixed
owner count. The lower bounds impose no anonymity, continuity,
measurability, homogeneity, or price representation on the original
mechanisms. The proof below constructs regular prices only as a consequence
of assuming a uniform bound for all finite owner models.

\subsection{Payment-only upper bound and a finite-owner lower bound}

For aggregate report $C=\sum_i b_i$, let $I=\{j:C_j>0\}$ and $k=|I|$.
For $j\in I$ define the independent-uniform face average
\begin{equation}\label{dim:uniform-rule}
 A_j(C)=\int_{[0,1]^{I\setminus\{j\}}}
 F\left(C_j e_j+\sum_{\ell\in I\setminus\{j\}}u_\ell C_\ell e_\ell\right)
 \,du_{-j},\qquad
 y_i(b)=\sum_{j\in I}\frac{b_{ij}}{C_j}A_j(C).
\end{equation}
An integral over an empty set of coordinates is its integrand; all payments
are zero when $I$ is empty.

\begin{proposition}\label{dim:uniform-upper}
The rule \eqref{dim:uniform-rule} is nonnegative, covers every coalition,
is TCWM at all real reports, and has total payment at most
$(k+1)F(C)/2$ when $k\ge1$.
\end{proposition}

\begin{proof}
For coalition aggregate $a$ and fixed outsider aggregate $o$, its
coordinate-$j$ payment is
\[
 \frac{a_j}{a_j+o_j}A_j(a+o).
\]
Both nonnegative factors are nondecreasing under any coordinatewise
increase of $a$. If a new coordinate becomes active, its own contribution
is nonnegative. In each old face average, integrating over that newly
active coordinate cannot lower the old integrand, whose coordinate was
identically zero. This proves finite TCWM, including comparisons between
different support faces.

For coverage, work in the $k$-dimensional active face. A finite concave
function is locally Lipschitz and differentiable almost everywhere on its
relative interior. At such a point $x$, its gradient $g(x)$ is nonnegative,
and homogeneity gives $g(x)\cdot x=F(x)$. Moreover
\begin{equation}\label{dim:gradient-majorant}
 g(x)\cdot z\ge F(z)
 \quad\text{for every nonnegative $z$ supported on this face.}
\end{equation}
This includes boundary $z$: for $0<t<1$, apply the supporting inequality
at $x$ to $(1-t)x+tz$, and combine it with concavity along the same segment.
After subtracting $F(x)$ and dividing by $t$, Euler's identity gives
\eqref{dim:gradient-majorant}, without taking a boundary-continuity limit.

Let $U_j$ be independent uniform random variables. For almost every fixed
$U_{-j}$, one-dimensional absolute continuity on $[\varepsilon,1]$ gives
\[
 C_j\int_\varepsilon^1
 \partial_jF(u_jC_j,U_{-j}C_{-j})\,du_j
 =F(C_j,U_{-j}C_{-j})-F(\varepsilon C_j,U_{-j}C_{-j})
 \le F(C_j,U_{-j}C_{-j}).
\]
The derivatives are nonnegative. Letting $\varepsilon\downarrow0$ and
using Tonelli therefore yields
$\E g_j(U\mathbin{\odot}C)\le A_j(C)/C_j$.
Averaging \eqref{dim:gradient-majorant} proves
$\sum_jz_jA_j(C)/C_j\ge F(z)$. Taking $z$ to be the aggregate of any
actual coalition proves coverage.

Finally, two applications of concavity give
\begin{align*}
 \sum_{j\in I}A_j(C)
 &\le\sum_{j\in I}
 F\left(C_je_j+\tfrac12\sum_{\ell\in I\setminus\{j\}}C_\ell e_\ell\right)\\
 &\le kF\left(\frac{k+1}{2k}C\right)
 =\frac{k+1}{2}F(C).
\end{align*}
For $k=1$ the same formula pays exactly $F(C)$; the empty face is immediate.
\end{proof}

\begin{proposition}[A finite obstruction]\label{dim:finite-lower}
Let $d\ge2$ and $s\ge2$ be integers. On $n=(d-1)s+1$ fixed owners,
each eligible for every resource and each with endowment box $[0,1]^d$,
the technology $F(C)=\min_jC_j$ forces every nonnegative,
coalition-covering TCWM rule to have funding factor at least
\begin{equation}\label{dim:finite-bound}
 1+\frac{(s-1)((d-1)s-1)}{2s^2+(d-3)s-1}.
\end{equation}
\end{proposition}

\begin{proof}
Suppose the factor is $\kappa$, and put $\delta=\kappa-1\ge0$.
Average the rule over all permutations of the owner labels and resource
labels. The fixed domain and technology are invariant under these finite
groups. Nonnegativity, every coverage and TCWM inequality, and the funding
bound are linear and hence are preserved. It suffices to consider this
invariant rule; symmetry is not an assumption on the original mechanism.

At profile $P_0$, put $s$ pure unit reporters in each resource group
$1,\ldots,d-2$, one pure unit reporter in each of the last two groups,
and leave the other $s-1$ actual owners at zero reports. Let the respective
per-owner payments be $a,b,w$, where $w\ge0$. Selecting one owner from
each resource group produces one, and grand production is one. Thus
\[
 (d-2)a+2b\ge1,\qquad
 (d-2)sa+2b+(s-1)w\le\kappa.
\]
Subtracting gives $(d-2)(s-1)a+(s-1)w\le\delta$, whence
\begin{equation}\label{dim:b-bound}
 2b\ge1-\frac{\delta}{s-1}.
\end{equation}
When $d=2$ the terms involving the absent first groups are zero.

At $P_1$, the $s-1$ zero reporters restore a pure unit of resource $d-1$.
TCWM applied to these owners together with the original resource-$(d-1)$
owner shows that this group's new payment is at least $b+(s-1)w\ge b$.
The first $d-1$ groups are now symmetric. If $S$ is their union, then
$Y_S(P_1)\ge(d-1)b$.

At $P_2$, every member of this same coalition $S$ restores its report to
one in each of resources $1,\ldots,d-1$ and zero in resource $d$.
These are feasible increases within the original unit boxes. Leave the
last-resource owner unchanged. Write $m=(d-1)s$. All $m$ members of $S$
now have the same payment $p$, and write $q$ for the last owner's payment.
One member of $S$ together with that owner produces one, so
$p+q\ge1$ while $mp+q\le\kappa$. Consequently
\[
 (d-1)b\le Y_S(P_1)\le Y_S(P_2)=mp
 \le\frac{m}{m-1}\delta.
\]
Combining this with \eqref{dim:b-bound} and rearranging gives
\eqref{dim:finite-bound}. All denominators are positive. All $n$ owners
remain present throughout; in particular, the proof never assumes that a
zero reporter receives zero at a positive-output profile.
\end{proof}

Letting $s\to\infty$ in \eqref{dim:finite-bound} and using
Proposition~\ref{dim:uniform-upper} proves $D_d=(d+1)/2$ for $d\ge2$.
For $d=1$, homogeneity makes $F$ linear on its ray; linear payments and
grand-coalition coverage give $D_1=1$.

\subsection{The logarithmic rule for full outside production}

Put $t_0=1-e^{-1/c}$. The defining equation for $c$ gives
$c(1+t_0)=1$. The function $h(x)=x(2-e^{-1/x})$ has
\[
 h'(x)=2-e^{-1/x}(1+1/x)>1\quad(x>0),
\]
and $h(1/2)<1<h(1)$, proving the asserted existence and uniqueness.
Let $U$ have density
\begin{equation}\label{dim:log-law}
 f_U(u)=\frac{c}{1-u}\ind_{(0,t_0)}(u).
\end{equation}
It has mass $-c\log(1-t_0)=1$ and mean
$c[-t_0-\log(1-t_0)]=c$.
Use one common scalar within each face average:
\begin{equation}\label{dim:log-rule}
 H_j(C)=\E F\left(C_je_j+U\sum_{\ell\ne j}C_\ell e_\ell\right),
 \qquad y_i(b)=\sum_{j:C_j>0}\frac{b_{ij}}{C_j}H_j(C).
\end{equation}
Absent-coordinate terms are zero. The independent-uniform and common-scalar
averages are different rules.

\begin{lemma}[Two scalar inequalities]\label{dim:scalar}
For positive $x\in\R_{++}^k$, $k\ge2$,
\begin{equation}\label{dim:ray-lemma}
 \sum_{j=1}^k\frac{\E\min\{x_j,U\min_{\ell\ne j}x_\ell\}}{x_j}\ge1.
\end{equation}
If $H:[0,1]\to\R_+$ is nondecreasing and concave and $0\le t\le t_0$,
then, with derivatives taken almost everywhere,
\begin{equation}\label{dim:hinge-lemma}
 \int_0^{t_0}[H(u)-uH'(u)]\,du
 =2\int_0^{t_0}H(u)\,du-t_0H(t_0)
 \ge t_0H(t)-tH(1).
\end{equation}
\end{lemma}

\begin{proof}
For \eqref{dim:ray-lemma}, choose one smallest coordinate $m$, let $S$
be the second-smallest coordinate, counting multiplicity, and put $t=m/S$.
The smallest coordinate contributes
\[
 q(t)=\E\min\{1,U/t\}
 =\begin{cases}
 1-c-c(1-t)\log(1-t)/t,&0<t\le t_0,\\
 c/t,&t_0\le t\le1.
 \end{cases}
\]
A second-smallest coordinate contributes $ct$, and all other terms are
nonnegative. If $t\le t_0$, then
$q(t)+ct-1=c(1-t)[-1-\log(1-t)/t]\ge0$.
If $t\ge t_0$, then $q(t)+ct=c(t+t^{-1})\ge2c>1$.

For \eqref{dim:hinge-lemma}, a possible downward jump at zero may be
filled with the right limit. This does not change the integral and only
strengthens the assertion when $t=0$. A finite nondecreasing concave
function is continuous from the left at one: concavity applied to a segment
ending at one proves this by letting the segment parameter tend to one.
After filling zero, therefore, $H$ is continuous on $[0,1]$.
The integral identity follows by integration by parts from a positive
endpoint and then passage to zero; boundedness gives $uH(u)\to0$.

A continuous piecewise-linear nondecreasing concave function is a
nonnegative constant plus a nonnegative linear function and a positive
linear combination of hinges $\min\{u,s\}$, $0<s\le1$.
For a constant $A$, the residual left side minus right side of the
inequality is $tA$; for $Bu$ it is $tB(1-t_0)$.
For the hinge with parameter $s$, the residual is respectively
\[
 \begin{cases}
 s(t-s),&s\le t,\\
 (s-t)(t_0-s),&t<s<t_0,\\
 t(s-t_0),&s\ge t_0.
 \end{cases}
\]
Each is nonnegative. Linear interpolation on increasingly fine grids
approximates a continuous concave function uniformly by such concave
piecewise-linear functions. The integral identity, involving only values
of $H$, passes the inequality to the limit.
\end{proof}

\begin{proposition}\label{dim:out-upper}
The rule \eqref{dim:log-rule} satisfies full OUT and coalition coverage at
every real report, and its total payment is at most
$[1+c(k-1)]F(C)$ when $k\ge1$.
\end{proposition}

\begin{proof}
Nonnegativity is immediate. As in the proof of
Proposition~\ref{dim:uniform-upper}, the coalition's coordinate-$j$
payment is the product of its nondecreasing share $a_j/(a_j+o_j)$ and
the nondecreasing face value $H_j(a+o)$. The same formula also shows that
an activated coordinate cannot lower any old contribution. Thus finite
TCWM holds, including across support faces. Two Jensen inequalities give
\[
 \sum_jH_j(C)
 \le\sum_jF\left(C_je_j+c\sum_{\ell\ne j}C_\ell e_\ell\right)
 \le[1+c(k-1)]F(C).
\]

First suppose $F$ is smooth on the positive relative interior of every
face. Fix a positive aggregate $C$ in its active face, coalition holdings
$0\le a\le C$, and a nonnegative restoration direction $z$ on this face.
Write $\alpha_j=a_j/C_j$ and
$v_j(u)=C_je_j+u\sum_{\ell\ne j}C_\ell e_\ell$.
Differentiating while both $a$ and $C$ increase in direction $z$ gives
\begin{align}
 D Y_S[z]&=\sum_j\{(1-\alpha_j)L_j(z)+\alpha_j M_j(z)\},
 \label{dim:holding-derivative}\\
 L_j(z)&=\frac{H_j(C)}{C_j}z_j,\nonumber\\
 M_j(z)&=\E\left[U\nabla F(v_j(U))\cdot z
 +(1-U)\partial_jF(v_j(U))z_j\right].\nonumber
\end{align}
The gradients are nonnegative and satisfy
\eqref{dim:gradient-majorant}, so $L_j\ge0$ and $M_j\ge cF(z)$.
The expression is affine in each holding fraction $\alpha_j$.
It is therefore enough to prove $DY_S[z]\ge F(z)$ at every corner
$\alpha\in\{0,1\}^k$.

The result is immediate if $F(z)=0$. For $k=1$ the technology is linear
on the active ray, and the assertion follows directly; assume $k\ge2$.
Initially suppose $z$ is strictly
positive on the active face. If at least two holding fractions are one,
the derivative is at least $2cF(z)\ge F(z)$.
If all are zero, put $x_j=C_j/z_j$. Monotonicity and homogeneity imply
\[
 \frac{F(v_j(u))}{F(z)}
 \ge\min\{x_j,u\min_{\ell\ne j}x_\ell\}.
\]
Equation~\eqref{dim:ray-lemma} yields $\sum_jL_j(z)\ge F(z)$.

It remains to consider a corner with only $\alpha_j=1$.
Normalize $G(x)=F(z\mathbin{\odot}x)/F(z)$, so $G(\mathbf1)=1$.
Common rescaling of $x=C/z$ leaves the normalized directional expressions
in \eqref{dim:holding-derivative} unchanged. We may therefore arrange
$\min_{\ell\ne j}x_\ell=1$. Put $t=x_j$ and $V=G(x)$.
Every face vector is at least $u x$, so the other coordinate terms satisfy
\[
 \sum_{\ell\ne j}\frac{L_\ell(z)}{F(z)}
 \ge cV\sum_{\ell\ne j}\frac1{x_\ell}\ge cV.
\]
The first part of $M_j/F(z)$ is at least $c$.
If $V\ge t_0$, their sum is at least $c(1+t_0)=1$.
If $t\ge1$, monotonicity gives $V\ge1$, so this case is already covered.
The only remaining case is $0<t\le V<t_0$.

Set $H(u)=G(te_j+u\sum_{\ell\ne j}x_\ell e_\ell)$.
Then $H(t)\ge t$, $H(1)=V$, and Euler's identity gives
\[
 \partial_jG(te_j+u x_{-j})
 =\frac{H(u)-uH'(u)}{t}.
\]
Because $(1-u)f_U(u)=c$ on $(0,t_0)$, the second part of $M_j/F(z)$
is at least
\[
 \frac{c}{t}\int_0^{t_0}[H(u)-uH'(u)]\,du
 \ge c(t_0-V)
\]
by \eqref{dim:hinge-lemma}. The singleton corner therefore contributes at
least $c+c(t_0-V)+cV=1$. This completes every holding corner.

Differentiation under the expectation is justified locally by Euler's
identity and nonnegative gradients:
\[
 C_j\partial_jF(v_j(u))\le F(C),\qquad
 uC_\ell\partial_\ell F(v_j(u))\le F(C)\quad(\ell\ne j).
\]
These bounds dominate the differentiated integrands on every compact
positive capacity set, including as $u\downarrow0$.
For directions with zero coordinates, apply the proved inequality to
$z+\eta\mathbf1_I$ and let $\eta\downarrow0$. The derivative expression
is a finite linear function of its direction, while monotonicity gives
$F(z+\eta\mathbf1_I)\ge F(z)$. No continuity of $F$ across faces is used.

For a finite restoration $z$, take $a(s)=a(0)+sz$ and
$C(s)=o+a(s)$, $0\le s\le1$. For every $s>0$ its active face is the
fixed face of $C(1)$. The derivative inequality integrates on
$[\varepsilon,1]$ to
$Y_S(1)-Y_S(\varepsilon)\ge(1-\varepsilon)F(z)$.
Finite TCWM gives $Y_S(\varepsilon)\ge Y_S(0)$ even if the initial
activation produces a payment jump. Letting $\varepsilon\downarrow0$
proves full OUT for smooth $F$. A one-coordinate face is handled by its
linear formula and the zero face is immediate.

Finally, choose a smooth probability density for $Z$ supported in
$[-1,1]^d$ and let
\[
 F_\varepsilon(x)=\E_Z
 F\bigl(\operatorname{diag}(e^{\varepsilon Z_j})x\bigr).
\]
This preserves nonnegativity, monotonicity, concavity, and degree-one
homogeneity. Convolution in logarithmic coordinates makes it smooth on
every positive relative face. Monotonicity and homogeneity give the
pointwise bounds
\begin{equation}\label{dim:smoothing-bound}
 e^{-\varepsilon}F(x)\le F_\varepsilon(x)
 \le e^\varepsilon F(x)
 \qquad\text{on every face.}
\end{equation}
The smooth finite OUT inequalities thus pass to the limit, since
\eqref{dim:smoothing-bound} also bounds each face average and each
nonnegative owner payment multiplicatively. This proves the claim for the
entire stated class, including downward boundary jumps.
Restoring a coalition's entire bundle from its zero report then proves
coverage: its initial payment under \eqref{dim:log-rule} is zero.
\end{proof}

\subsection{Why arbitrary finite-owner rules satisfy the matching OUT lower bound}

The remaining argument concerns only the continuous technology
$F_r(C)=\min_{1\le j\le r}C_j$. Its continuity imposes no restriction on
the broader class used for the upper bound. We give the reduction in full
to distinguish a fixed finite-owner impossibility from a restriction to
price mechanisms.

\begin{lemma}[Prices forced by bounded finite-owner models]
\label{dim:price-reduction}
Fix $r\ge2$. Suppose that, for every $N$, the model with $N$ fixed
pure-resource owners of each of the $r$ types, each with report interval
$[0,1]$, admits a nonnegative, coalition-covering, full-OUT rule funded by
$\kappa F_r$. Then there are nonnegative degree-zero homogeneous prices $p(C)$ on
$\R_{++}^r$, locally Lipschitz there, such that
\begin{align}
 C\cdot p(C)&\le\kappa F_r(C),\label{dim:price-budget}\\
 a\cdot p(C)&\ge F_r(a)\qquad(0\le a\le C),
 \label{dim:price-coverage}\\
 a\cdot p(C)&\ge(a-u)\cdot p(C-u)+F_r(u)
 \qquad(0\le u\le a\le C).
 \label{dim:price-out}
\end{align}
When an aggregate has a zero coordinate, all its prices may be set to
zero and these statements still hold at every nonnegative real profile.
The prices can also be chosen symmetric under resource permutations.
\end{lemma}

\begin{proof}
For each $N$, average its rule over permutations within each resource
type. At a binary profile with $m_j$ active unit reporters of type $j$,
write $\pi_j^{(N)}(m)$ for an active owner's payment. For integer
$0\le u\le a\le m$, coverage, budget, and the OUT comparison that
turns $u_j$ coalition members of type $j$ completely off give
\begin{align}
 m\cdot\pi^{(N)}(m)&\le\kappa F_r(m),\nonumber\\
 a\cdot\pi^{(N)}(m)&\ge F_r(a),\label{dim:count-array}\\
 a\cdot\pi^{(N)}(m)&\ge
 (a-u)\cdot\pi^{(N)}(m-u)+F_r(u).\nonumber
\end{align}
Inactive reporters may have positive payments. We have only dropped their
nonnegative payments from the budget and from the right side of the last
inequality. Set an unused coordinate $\pi_j(m)$ to zero when $m_j=0$.
For $m_j>0$ the coordinate bound is
$0\le\pi_j^{(N)}(m)\le\kappa F_r(m)/m_j$, independently of $N$.
A diagonal subsequence as $N\to\infty$ therefore gives one array
$\pi(m)$ satisfying \eqref{dim:count-array} for all integer count vectors.
Every finite group of these constraints is available once $N$ exceeds
its largest count. The set $\mathcal C$ of arrays obeying these constraints
and coordinate bounds is nonempty, compact in the countable product
topology, and convex.

For a positive integer $h$ let $(T_h\pi)(m)=\pi(hm)$.
Multiplying all counts in \eqref{dim:count-array} by $h$ and dividing by
$h$ shows that $T_h$ preserves $\mathcal C$. These linear maps commute.
For the first $L$ primes $\ell_1,\ldots,\ell_L$, average $T_h\pi$ over
$h=\prod_{j=1}^L\ell_j^{e_j}$, $0\le e_j<L$.
For a fixed prime $\ell_j$, shifting its exponent changes this average
by at most $\kappa F_r(m)/(Lm_v)$ in coordinate $(m,v)$, by cancellation
of the interior terms. A diagonal subsequential limit as $L\to\infty$
is fixed by every prime dilation, hence by every positive integer dilation.

For rational $C>0$ define $p(C)=\pi(LC)$, where $L$ clears the
denominators. Integer-dilation invariance makes this independent of $L$
and homogeneous under positive rational scaling. Clearing the
denominators of $C,a,u$ in \eqref{dim:count-array} proves
\eqref{dim:price-budget}--\eqref{dim:price-out} for rational data.

Regularity follows from finite comparisons, rather than from any property
of the original rules. For rational $C,C'>0$, put
$\lambda=\min_j C_j/C'_j$ and $v=C-\lambda C'\ge0$.
At aggregate $\lambda C'$, give a coalition all of resource $j$ and no
other resource; then restore $v$ to that same coalition. TCWM, which is
implied by OUT, and rational homogeneity yield
\begin{equation}\label{dim:relative-comparison}
 p_j(C')-p_j(C)
 \le\frac{v\cdot p(C)}{\lambda C'_j}.
\end{equation}
On any compact box $[a,b]^r$ with $a>0$, the coordinate prices are bounded
by $\kappa b/a$, the denominator is at least $a^2/b$, and
\[
 0\le v_\ell\le\|C-C'\|_\infty(1+b/a).
\]
Interchanging $C,C'$ supplies the opposite inequality. Thus $p$ is
Lipschitz on the rational points of each such box and has a unique locally
Lipschitz extension to the positive orthant. Its inequalities and
homogeneity pass to real data by continuity. If coalition holdings or
withheld amounts have zero coordinates, approximate retained holdings $a-u$, withdrawals $u$, and outsider
holdings $C-a$ separately by nonnegative rational vectors, retaining their zero
patterns, and reconstruct $a,C$ by summation. This preserves
$0\le u\le a\le C$; both positive aggregate totals remain positive.

If $F_r(C)=0$, every subbundle of $C$ has zero worth, so zero prices give
coverage and budget on that face. If $C>0$ but $F_r(C-u)=0$, the right
payment in \eqref{dim:price-out} is zero, while coverage and monotonicity
give $a\cdot p(C)\ge F_r(a)\ge F_r(u)$.
These observations complete all boundary comparisons. Finally, average
the constructed prices over the finite resource-permutation group.
All the inequalities, homogeneity, and local Lipschitz continuity survive,
and the resulting prices are resource-symmetric.
\end{proof}

\begin{lemma}[Removing abundant coordinates]\label{dim:pair-limit}
Under the hypothesis of Lemma~\ref{dim:price-reduction}, there is a
symmetric, homogeneous, locally Lipschitz price rule for $\min\{A,B\}$
with nonnegative payments, coalition coverage, full OUT, finite funding
factor $\kappa$, and the expenditure cap
\begin{equation}\label{dim:pair-cap}
 B P_2(A,B)\le D A\quad(B\ge A>0),
 \qquad D=\frac{\kappa-1}{r-1},
\end{equation}
together with its symmetric counterpart.
\end{lemma}

\begin{proof}
At $(A,M,\ldots,M)$, $0<A<M$, symmetry gives a short-resource price
$q$ and common long-resource price $b$. Coverage of a bundle with $A$
units of every resource and the budget yield
\[
 q+(r-1)b\ge1,\qquad Aq+(r-1)Mb\le\kappa A.
\]
Subtracting $A$ times the first inequality gives
$Mb\le D A M/(M-A)$.
At $(A,B,M,\ldots,M)$ with $A\le B\le M$, restore a coalition owning
all $B$ units of the second resource to $M$ units. TCWM gives
\begin{equation}\label{dim:prelimit-cap}
 Bp_2(A,B,M,\ldots,M)\le D\frac{AM}{M-A}.
\end{equation}
For $r=2$, the left side is independent of $M$; letting $M\to\infty$
already proves the lemma using the prices from
Lemma~\ref{dim:price-reduction}. Suppose henceforth $r\ge3$.

For each positive rational pair $(A,B)$, the first two prices at
$(A,B,M,\ldots,M)$ lie in $[0,\kappa]$. A common diagonal subsequence
$M\to\infty$ gives rational pair prices $P_1,P_2$.
They inherit symmetry and pair budget at most $\kappa\min\{A,B\}$.
Equation~\eqref{dim:prelimit-cap} gives \eqref{dim:pair-cap}.
For coalition pair holdings $(a,b)\le(A,B)$, add $\min\{a,b\}$ units
of every extra resource to this coalition. The added payment is at most
\[
 \min\{a,b\}\sum_{j\ge3}p_j(A,B,M,\ldots,M)
 \le\frac{\kappa\min\{a,b\}\min\{A,B\}}{M},
\]
which vanishes. Its worth is $\min\{a,b\}$, so pair coverage survives.

For OUT, it is essential to control the shifted extra totals in a
deviation. For fixed $w\ge0$, compare
$z=(A,B,M,\ldots,M)$ and $z'=(A,B,M-w,\ldots,M-w)$, where $M>w$.
Equation~\eqref{dim:relative-comparison}, first in one orientation and
then the other, gives for $j=1,2$
\begin{align}
 p_j(z')-p_j(z)&\le\frac{w\kappa\min\{A,B\}}{Mz_j},
 \label{dim:shift-bound}\\
 p_j(z)-p_j(z')&\le
 \frac{w\kappa\min\{A,B\}}{(M-w)z_j}.\nonumber
\end{align}
Indeed, the forward residual consists only of $w$ units in each extra
coordinate; the reverse scaling is $(M-w)/M$ and its residual lies in
the first two coordinates. The budget bounds both residual expenditures.

Fix rational pair holdings $(a,b)$ and withdrawals $(u,v)\le(a,b)$,
and put $w=\min\{u,v\}$. Give that coalition $w$ units of every extra
resource and withdraw all of them. Its full outside production is $w$.
The high extra-coordinate payment vanishes as above, and its low retained
extra holdings are zero. At the low pair $(A-u,B-v)$ with positive
production, \eqref{dim:shift-bound} shows that the prices with extra totals
$M-w$ have the same subsequential limit as those with extra totals $M$.
If that low pair has zero production, its entire payment is zero by budget.
The limiting inequality is therefore
\begin{align*}
 aP_1(A,B)+bP_2(A,B)
 \ge{}&(a-u)P_1(A-u,B-v)\\
 &+(b-v)P_2(A-u,B-v)+\min\{u,v\}.
\end{align*}
Set prices to zero at zero-production pairs. We have obtained a rational
pair rule satisfying all required finite inequalities and the cap.

This limit need not yet be homogeneous. Consider the compact convex set
of rational pair-price arrays obeying these inequalities, symmetry, the
budget, and \eqref{dim:pair-cap}. Each coordinate is bounded by
$\kappa\min\{A,B\}/A$ or $\kappa\min\{A,B\}/B$.
The operations $(T_hP)(A,B)=P(hA,hB)$ preserve the set. Repeat the
prime-dilation averaging in Lemma~\ref{dim:price-reduction}; all its
constraints are homogeneous and closed, so the resulting rational array
is homogeneous and retains the cap. The finite comparison
\eqref{dim:relative-comparison} now gives local Lipschitz continuity and
extension to real positive pairs exactly as before. Coverage handles any
zero-production low profile as in the last paragraph of that lemma.
\end{proof}

\begin{lemma}[The two-resource obstruction]\label{dim:pair-obstruction}
Every pair-price rule with the properties of
Lemma~\ref{dim:pair-limit} has $D\ge c$.
\end{lemma}

\begin{proof}
For $A\ge B>0$ write $t=B/A$, $p(t)=P_1(A,B)$ and $q(t)=P_2(A,B)$.
These are well defined by homogeneity, locally Lipschitz for $0<t\le1$,
and $p(t)\le Dt$. Budget gives $q(t)\le\kappa$, while symmetry at the
positive balanced profile gives $q(1)=p(1)\le D$.

If a coalition owns all of the short resource and restores only that
resource, TCWM implies that $z(t)=tq(t)$ is nondecreasing. To obtain the
OUT constraint, start at aggregate $(A,B)$ with coalition holdings
$(0,B)$ and restore $(h,h)$ to that same coalition. Its new holdings are
$(h,B+h)$ and its new aggregate is $(A+h,B+h)$, with outsiders fixed.
This is a feasible restoration from a zero holding, not a withdrawal
from that zero holding. At almost every $t<1$, differentiating its finite
OUT inequality at $h=0+$ gives
\[
 p(t)+q(t)+t(1-t)q'(t)\ge1.
\]
Using the cap, this yields
\[
 \left[\frac{tq(t)}{1-t}\right]'
 =\frac{q(t)+t(1-t)q'(t)}{(1-t)^2}
 \ge\frac{1-Dt}{(1-t)^2}.
\]
Integrate from $\varepsilon$ to $t$ and send $\varepsilon\downarrow0$.
Boundedness of $q$ makes the lower endpoint vanish, and hence
\[
 z(t)\ge (1-D)t-D(1-t)\log(1-t).
\]
Since $z$ is nondecreasing and $z(1)=q(1)\le D$, the right side is at
most $D$ for every $0<t<1$. If $D=0$, this is already a contradiction.
For $D>0$, its maximum occurs at $t=1-e^{-1/D}$ and equals
$1-D+De^{-1/D}$. Consequently
$D(2-e^{-1/D})\ge1$. The strict monotonicity of $h$ proved above implies
$D\ge c$.
\end{proof}

\begin{proof}[Completion of Theorem~\ref{dim:main}]
The TCWM equality was proved after Proposition~\ref{dim:finite-lower}.
Proposition~\ref{dim:out-upper} proves the OUT upper bound.
For $d\ge2$, suppose a factor $\kappa<1+c(d-1)$ were feasible for
every fixed finite pure-owner unit-box model of $F_d=\min_jC_j$.
Lemmas~\ref{dim:price-reduction}--\ref{dim:pair-obstruction} would give
$(\kappa-1)/(d-1)\ge c$, a contradiction. Therefore for each such
$\kappa$ at least one finite, fixed, bounded owner model has no
$\kappa$-funded rule. This is a compactness contrapositive over finite
models, not a mechanism on an infinite owner population.
If the supremum of their infimal feasible factors were smaller than the
claimed constant, choosing $\kappa$ strictly between them would give just
such a rule for every model. This also deals with possible nonattainment
of individual infima. For $d=1$, linear payments satisfy full OUT with
factor one and coverage forces that factor.
\end{proof}

\begin{remark}\label{dim:scope}
These dimension frontiers concern protection of total coalition payments,
with the full value of withheld production added only in the OUT contract.
They do not settle a sharp frontier for individual payment monotonicity.
The displayed integrals are mathematical payment formulas; no efficient
evaluation guarantee for an arbitrary value oracle is asserted. Intermediate
outside-protection parameters and extra deletion-marginal floors are not
included in the present theorem.
\end{remark}

\section{Sparse production: a uniform coalition-size upper bound}\label{sec:upper}
We first state the upper bound with which the new lower must be compared.
The binomial semivalue and bounded-rank support inequalities are established
ingredients \citep{FeigeEtAl2014,EzraEtAl2018,Casajus2011}. Their use here retains
actual mixed owners, all-coalition coverage and the stronger full-OUT contract.
The proof is included to expose the precise comparison class.

\begin{theorem}[Sparse-production guarantee]\label{thm:sparse-upper}
For every finite packing institution of activity arity at most $r\ge1$ and
every integer $s\ge1$, there is an all-profile rule satisfying coalition
coverage and full OUT for every coalition of size at most $s$, with
\begin{equation}\label{eq:sparse-upper}
 K_r(s)\le K_r^{\mathrm{OUT}}(s)\le r s^{r-1}.
\end{equation}
For $p=1/s$, one such rule is
\begin{equation}\label{eq:semivalue}
 m_i(p;b)=\mathbb E_{R\subseteq N\setminus\{i\}}
       [F(b_i+b_R)-F(b_R)],\qquad y_i(b)=p^{1-r}m_i(p;b),
\end{equation}
where other owners are independently included in $R$ with probability $p$.
For rational data the rule is a finite rational computation. No
polynomial-time implementation or incentive guarantee for a sampled estimator
is asserted.
\end{theorem}

\begin{lemma}[Owner support and thinning]\label{lem:owner-support}
Fix a report profile and let $v(T)=F(b_T)$. For $i\in R$, set
$U_i(R)=v(R)-v(R\setminus\{i\})$. Then
\begin{align}
 \sum_{i\in T}U_i(R)&\ge v(T) &&(T\subseteq R),\label{eq:dual-cover}\\
 \sum_{i\in R}U_i(R)&\le r v(R),\label{eq:rank-loss}\\
 \mathbb E v(S\cap R_p)&\ge p^r v(S),\label{eq:thinning}
\end{align}
where $R_p$ independently retains every owner with probability $p$.
\end{lemma}
\begin{proof}
An optimal nonnegative dual price $w$ at $b_R$ satisfies
$v(R)=w\cdot b_R$ and $v(R\setminus\{i\})\le w\cdot b_{R\setminus\{i\}}$.
Thus $U_i(R)\ge w\cdot b_i$, and dual feasibility proves
\eqref{eq:dual-cover}.

For a fixed $S$, choose an optimal production plan $x$. For every used
activity $e$ and each required row $a$, assign that row to an actual supplier
$i\in S$ with probability $b_{ia}/(b_S)_a$, independently across required
rows. The denominator is positive whenever $x_e>0$ and $A_{ae}>0$.
Each assignment gives a nonnegative value piece, of weight
$\pi_e x_e$ times its assignment probability, supported on at most $r$ owners.
Let $h(T)$ be the total weight of pieces whose suppliers all belong to
$T\subseteq S$. For activity $e$, its retained fraction is
\[
 \prod_{a:A_{ae}>0}\frac{(b_T)_a}{(b_S)_a},
\]
which is at most the fraction in each individual required row. The retained
plan is therefore feasible for $b_T$. Consequently $h(T)\le v(T)$ and
$h(S)=v(S)$. Applying this at $S=R$ gives
\[
 U_i(R)\le h(R)-h(R\setminus\{i\});
\]
summing counts each value piece at most $r$ times and proves
\eqref{eq:rank-loss}. Under independent thinning, a piece supported on $k\le r$
owners survives with probability $p^k\ge p^r$. Hence
$\mathbb E v(S\cap R_p)\ge\mathbb E h(S\cap R_p)\ge p^r v(S)$.
This decomposes a production plan, not the strategic owners.
\end{proof}

\begin{proof}[Proof of Theorem~\ref{thm:sparse-upper}]
Let $R_p$ now thin all owners. Conditioning on whether $i\in R_p$ gives
\[
 p\sum_{i\in S}m_i(p;b)
 =\mathbb E\sum_{i\in S\cap R_p}U_i(R_p)
 \ge\mathbb E v(S\cap R_p)\ge p^r v(S).
\]
After multiplication by $p^{-r}$, this is coverage of every $S$, including
$|S|>s$. Nonnegativity follows from monotonicity. Also
\[
 p\sum_i m_i(p;b)\le r\mathbb E F(b_{R_p})
 \le rF(pB)=rpF(B),
\]
where the second inequality is concavity. This proves the funding bound.

For $s\ge2$, fix a protected coalition $C$ and write $M_C=\sum_{i\in C}m_i$.
Expanding the finite expectation yields
\[
 M_C(b)=\sum_{A\subseteq N}c_C(A)F(b_A),
\]
with, for $0<|A|<n$,
\begin{equation}\label{eq:coefficient}
 c_C(A)=p^{|A|-1}(1-p)^{n-|A|-1}(|A\cap C|-p|C|).
\end{equation}
The endpoints are $c_C(\varnothing)=-|C|(1-p)^{n-1}$ and
$c_C(N)=|C|p^{n-1}$. Since $p|C|\le1$, every coefficient of a value affected
by $C$'s report is nonnegative. For a withdrawal $\delta$ supported on $C$,
superadditivity gives
$F(b_A)-F(b_A-\delta_A)\ge F(\delta_A)$.
For $A\cap C=\varnothing$, both this difference and $F(\delta_A)$ vanish.
Therefore
\[
 M_C(b)-M_C(b-\delta)\ge\sum_A c_C(A)F(\delta_A)=M_C(\delta).
\]
The profile $\delta$ lies in the same endowment boxes, with the same owners
and the same rule. Coverage at that profile gives
$p^{1-r}M_C(\delta)\ge F(\delta_C)$, proving full OUT.

At $s=1$, $p=1$ and $m_i=F(B)-F(B-b_i)$. The second term is fixed under
$i$'s withdrawal, so superadditivity proves full OUT directly; the first two
inequalities of Lemma~\ref{lem:owner-support} give coverage and factor $r$.
If no positive-profit activity exists, the zero rule suffices. Empty
coalitions and the empty owner set are immediate.
\end{proof}

\section{A square-root lower bound for arity three}\label{sec:lower}
The following lower applies to arbitrary rules under the full report-domain
contract. It does not impose a marginal formula, an efficient-price floor,
anonymity, or a zero payment to inactive owners.

\begin{theorem}[Projective cap exchange]\label{thm:projective}
For each integer $M\ge1$, there is a unit-profit packing technology of arity
three on $d=2^{M+1}-1$ resources, with $n=2M+1$ fixed fully eligible owners,
for which every coalition-covering rule satisfying TCWM for coalitions of
size at most $2M-1$ has funding factor at least
\begin{equation}\label{eq:projective-finite}
 \max\left\{\frac23 Q_M,\ 1+\frac{2(M-1)}{3M}Q_M\right\},
 \qquad Q_M=\prod_{m=2}^M\frac{2m}{2m-1}
             =\frac{4^M}{2\binom{2M}{M}}.
\end{equation}
Every owner's report domain is the full real unit cube. Consequently, for
every integer $s\ge1$,
\begin{equation}\label{eq:main-interval}
 \sqrt{s/6}\le K_3(s)\le K_3^{\mathrm{OUT}}(s)\le3s^2.
\end{equation}
\end{theorem}

\subsection{Technology, legal symmetry and productive profiles}
Let $V=\mathbb F_2^{M+1}$ have basis $v,e_1,\ldots,e_M$.
Resources are the nonzero vectors of $V$. For every unordered triple
$\{a,b,a+b\}$ of distinct nonzero vectors, include one activity consuming one
unit of each row and yielding one unit of output. Duplicate triples are
identified. All activities have exactly three inputs.

There are actual owners $p,h_1,\ell_1,\ldots,h_M,\ell_M$, each endowed with
the all-ones vector. Every report in $[0,1]^d$ is legal for every owner.
The product of owner permutations and invertible binary linear maps acts on
this institution: the latter permute resources and activities, while owner
permutations preserve the identical report domains.

\begin{lemma}[Finite averaging]\label{lem:average}
If a rule on this institution satisfies coverage, size-$s$ TCWM (or full OUT)
and factor $\kappa$, then its finite group average satisfies the same
constraints and factor, and is equivariant. A lower on total payment of the
average at a profile of production one is a lower on the original rule at
some profile in that profile's finite symmetry orbit.
\end{lemma}
\begin{proof}
For a group element $g$ with owner permutation $\sigma_g$, define the
conjugate rule by $y^g(b)=\sigma_g^{-1}y(gb)$. Resource and owner maps preserve
coordinate order, report feasibility, coalition values, coalition sizes,
and the production value of a withheld bundle. Each contract inequality
therefore holds for every conjugate. All constraints are linear in payment
values, so their finite average also satisfies them. A change of variable
in the group sum proves equivariance. Finally, its total payment at $b$ is
the average of original total payments at the legal profiles $gb$, which
all have production $F(B)$. At least one is at least that average. This is a
proof transformation of rules; no strategic comparison changes an owner.
\end{proof}

For $m=1,\ldots,M$, define profile $C_m$ as follows. Owner $p$ reports the
unit vector in row $v$; each $h_j$ reports the unit vector in row $e_j$;
$\ell_j$ reports the unit vector in row $v+e_j$ if $j\le m$, and zero
otherwise. Every unmentioned resource coordinate is zero. In particular,
\emph{dormant high owners still report $e_j$} when $j>m$.

\begin{lemma}[Exact support and value]\label{lem:star-support}
The only activities supported at $C_m$ are the $m$ triples
$\{v,e_j,v+e_j\}$ with $j\le m$, and $F(C_m)=1$.
\end{lemma}
\begin{proof}
In the possible row set
$\{v\}\cup\{e_j,v+e_j:1\le j\le M\}$, the sum of private rows from
 distinct indices has two nonzero $e$ coordinates and leaves the set. The
sum of the private pair at index $j$ is $v$, and adding $v$ exchanges those
two rows. Thus every supported triple is precisely a complete active branch.
Every branch consumes the reported unit of $v$, giving the upper bound one;
one branch attains it. Equivalently, the full packing dual prices $v$ and
all zero-capacity rows at one and all other rows at zero. Supported columns
contain $v$ and unsupported columns contain a zero-capacity row, so this is
a feasible dual of value one for the full technology.
\end{proof}

\subsection{A constant seed and a fixed-owner payment ladder}
Work with the averaged rule of Lemma~\ref{lem:average}. At $C_1$, the linear
map exchanging $v$ and $e_1$ and fixing $e_j$ for $j>1$, together with the
owner swap $(p,h_1)$, preserves the whole profile. So does the invertible map
$v\mapsto v+e_1$, fixing every $e_j$, together with the swap $(p,\ell_1)$.
These maps fix the dormant high reports. Hence the three productive owners
receive a common payment $a$. Their coalition produces one, and coverage
requires $3a\ge1$. This uses static coverage for a three-member coalition,
not a three-member incentive constraint.

At $C_m$, permuting active indices and independently replacing
$e_j$ by $e_j+v$ for active $j$ preserves the entire profile after the matching
owner swaps. These invertible maps fix $v$ and all inactive high rows.
They act transitively on the $2m$ active private owners. Let $q_m$ be their
common payment and put $T_m=2m q_m$. The seed is
\begin{equation}\label{eq:seed}
 q_1\ge\frac13,\qquad T_1\ge\frac23.
\end{equation}

The transition $C_{m-1}\to C_m$ restores only the actual owner $\ell_m$.
Protect the coalition consisting of $\ell_m$ and the $2m-2$ old active
private owners. Its cardinality is $2m-1$, and every old member's report
is unchanged. At the old profile its payment is at least $T_{m-1}$:
$y_{\ell_m}(C_{m-1})$ is nonnegative, and is \emph{not} assumed to be zero.
At the new profile the protected sum is $(2m-1)q_m$. Therefore TCWM gives
\begin{equation}\label{eq:ladder}
 (2m-1)q_m\ge2(m-1)q_{m-1},\qquad
 T_m\ge\frac{2m}{2m-1}T_{m-1}.
\end{equation}
The unchanged new high owner $h_m$ and the cap owner $p$ are not in this
protected coalition. No owner is split, cloned, added, or removed.
Iterating \eqref{eq:ladder} proves $T_M\ge2Q_M/3$. Since production is one,
nonnegativity and funding imply $\kappa\ge2Q_M/3$.

Let $z=y_p(C_M)$. Terminal branch coverage gives $z+2q_M\ge1$, and total
payment is at least $z+2M q_M$. Hence
\[
 \kappa\ge z+2M q_M\ge1+2(M-1)q_M
 \ge1+\frac{2(M-1)}{3M}Q_M.
\]
This remains valid even if $2q_M>1$; it does not set $z=1-2q_M$.
For an asymmetric original rule, Lemma~\ref{lem:average} locates the costly
profile somewhere in the orbit, not necessarily at the displayed $C_M$.
This proves the finite assertion of Theorem~\ref{thm:projective}.

\subsection{The uniform exponent and an optimized tail}
The elementary estimate
\begin{equation}\label{eq:central-bound}
 Q_M\ge\frac12\sqrt{3M+1}
\end{equation}
follows by induction. The base $M=1$ is equality; after squaring, the step
reduces to
$4(M+1)^2(3M+1)-(2M+1)^2(3M+4)=M\ge0$.
For $M=\lfloor(s+1)/2\rfloor$, all incentive coalitions used have size at
most $s$, and $M\ge s/2$. Thus
$K_3(s)\ge2Q_M/3\ge\sqrt{3M+1}/3\ge\sqrt{s/6}$.
Combining this with Theorem~\ref{thm:sparse-upper} proves
\eqref{eq:main-interval}, including $s=1$.

\begin{corollary}[Optimized ladder certificate]\label{cor:tail}
Set $B_1=1$. For $s\ge2$, let $h=\lceil s/2\rceil$ and
\begin{equation}\label{eq:Bs}
 B_s=1+\frac{2(s-1)Q_h}{3h}
            \left(\frac{s-1}{s}\right)^{s-h}.
\end{equation}
Then $K_3(s)\ge B_s$, and
\begin{equation}\label{eq:tail-asymptotic}
 \lim_{s\to\infty}\frac{B_s}{\sqrt s}
   =\frac{\sqrt{2\pi}}{3\sqrt e}.
\end{equation}
The witness may use $2s+1$ fully eligible owners and $2^{s+1}-1$ resource
rows. In particular $B_6=493/243>2$, with 13 owners and 127 rows.
The optimization is of this ladder certificate, not of the unrestricted
funding frontier.
\end{corollary}
\begin{proof}
Keep $s$ fixed and permit arbitrary ladder length $M$. At stage $m$ protect
$\ell_m$ and any $k_m=\min\{s-1,2m-2\}$ old active private owners. Whole
pairs are unnecessary; all selected old owners have payments $q_{m-1}$ and
$q_m$ at the respective profiles. TCWM and nonnegativity imply
\[
 (k_m+1)q_m\ge k_m q_{m-1},\qquad
 q_M\ge\tfrac13 P_{M,s},\qquad
 P_{M,s}=\prod_{m=2}^M\frac{k_m}{k_m+1}.
\]
Terminal coverage therefore gives
$\kappa\ge1+2(M-1)P_{M,s}/3$.
Up to $h$ the uncapped recurrence gives $P_{h,s}=Q_h/h$; afterwards every
factor is $\alpha=(s-1)/s$. For $M\ge h$ and $M>1$, the consecutive ratio
of $(M-1)P_{M,s}$ is $M\alpha/(M-1)$, which is greater than one for $M<s$,
equal to one for $M=s$, and less than one for $M>s$.
Before the cap the increment increases; its $M=1$ value is zero and is
handled separately. Thus the equal maxima occur at $M=s,s+1$, proving
\eqref{eq:Bs}. The central-binomial asymptotic
$Q_h\sim\sqrt{\pi h}/2$, together with $h\sim s/2$ and
$\alpha^{s-h}\to e^{-1/2}$, proves \eqref{eq:tail-asymptotic}.
Finally $Q_3=8/5$ gives $B_6=493/243$.
\end{proof}

\subsection{Exact certificate form and what computation verifies}
The displayed argument has nonnegative rational multipliers. For fixed $M,s$,
use the rows
\[
 3q_1\ge1,\qquad (k_m+1)q_m-k_mq_{m-1}\ge0,
 \qquad z+2q_M\ge1,\qquad \kappa-z-2M q_M\ge0.
\]
Give the two final rows weight one. Starting with residual coefficient
$2(M-1)$ at $q_M$, eliminate backwards: weight the row at $m$ by the
residual divided by $k_m+1$, then replace the residual by that weight times
$k_m$. Weight the seed by the final residual divided by three.
All coefficients of $q_m$ and $z$ cancel, leaving
$\kappa\ge1+2(M-1)P_{M,s}/3$. At $M=1$ only the final rows are needed.
A lifted certificate can expand averaging back into the original
owner/profile variables; it retains the nonnegativity rows for dormant
payments rather than replacing those payments by zero.

The accompanying standard-library checks verify small full projective
technologies, the required linear maps, exact finite coefficient cancellation,
and the optimized scalar tail. A selected original-variable lift additionally
checks 34,440 inequalities cancelling 26,040 payment coordinates. These are
reproducible finite certificates supporting the written argument. The universal
group averaging and all-size geometry are conventional proofs, not statements
proved merely by the finite tests.


\section{Full outside protection on a four-vertex path}
\label{graph:section}

This section gives a structured unit-profit graph guarantee. For a finite
resource graph $G$, let
\begin{equation}\label{graph:production}
 F_G(c)=\max\left\{\sum_{e\in E(G)}x_e:
 x_e\ge0,\quad\sum_{e\ni v}x_e\le c_v\ (v\in V(G))\right\}.
\end{equation}
Every edge consumes one unit of each endpoint resource and produces one
unit of output. Owners and their finite coordinatewise endowment boxes
remain fixed; all real reports in those boxes are allowed. Owners may
hold arbitrary mixtures of resources, and every coalition, including its
unchanged members, receives coverage and full outside-production protection.

\begin{theorem}[Protected-leaf path rule]\label{graph:theorem}
For the path $B-A-D-C$, there is an explicit nonnegative payment rule that
satisfies coalition coverage and full OUT for every coalition at every
real report, with
\begin{equation}\label{graph:funding-bound}
 \sum_i y_i(b)\le(1+1/\sqrt2)F_G\left(\sum_i b_i\right).
\end{equation}
The same guarantee holds for complete blow-ups of this path, their finite
disjoint unions, and additional isolated resources, with the same fixed
mixed owners. Thus the total-funding factor is $1+1/\sqrt2$, and the
extra subsidy is at most $F_G(\sum_i b_i)/\sqrt2$.
\end{theorem}

A complete blow-up replaces each of the four vertices by a nonempty
independent group, includes every edge between groups corresponding to a
path edge, and includes no other edges. The theorem is restricted to unit
profits and these graph classes. It asserts neither an all-graph bound
below two nor optimality among payment rules on the path.

\subsection{The selected flow and its payment rule}

Write the path capacities as $c=(A,B,C,D)$, using a vertex label also for
its capacity, and put
\[
 \alpha=\frac1{\sqrt2},\qquad \beta=\alpha^2=\frac12.
\]
These are the constants of this section. For the path,
\begin{equation}\label{graph:path-value}
 F_G(c)=\min\{A+C,B+D,A+D\}.
\end{equation}
Indeed, each expression is a vertex-cover capacity and hence an upper
bound. If $P$ is their minimum, the edge quantities
\[
 x_{AB}=(P-D)_+,\qquad x_{DC}=(P-A)_+,\qquad
 x_{AD}=P-(P-D)_+-(P-A)_+
\]
are feasible and sum to $P$. Here $t_+=\max\{t,0\}$;
nonnegativity of $x_{AD}$ follows from $P\le A+D$, and the other capacity
constraints follow from $P\le A+C$ and $P\le B+D$.

Form a directed bipartite double cover: a left copy of each vertex has
capacity $c_v$, a right copy has capacity $\alpha c_v$, and an edge $uv$
gives arcs $u_Lv_R$ and $v_Lu_R$. Its two components have maximum values
\begin{equation}\label{graph:oriented-values}
 V=\min\{A+C,\alpha(B+D),A+\alpha D\},\qquad
 V'=\min\{B+D,\alpha(A+C),D+\alpha A\}.
\end{equation}
We require a particular choice of left uses, rather than an arbitrary
maximum flow.

\begin{lemma}[All maximum-flow left uses]\label{graph:flow-interval}
In the component $\{A_L,C_L\}\to\{B_R,D_R\}$, the possible left uses
of a maximum flow are exactly
\begin{equation}\label{graph:first-interval}
 w_A+w_C=V,\qquad
 \max\{0,V-A\}\le w_C\le\min\{C,\alpha D,V\}.
\end{equation}
In the other component they are exactly
\begin{equation}\label{graph:second-interval}
 w_B+w_D=V',\qquad
 \max\{0,V'-D\}\le w_B\le\min\{B,\alpha A,V'\}.
\end{equation}
Both intervals are nonempty, including at zero capacities.
\end{lemma}
\begin{proof}
The three bounds defining $V$ are the total left capacity, total right
capacity, and the combined capacity of $A_L,D_R$. They upper-bound every
flow. For any value-$V$ flow, $w_C\le C$, $w_C\le\alpha D$ because
$C_L$ can only send to $D_R$, and $0\le w_C\le V$.
The bound $w_A\le A$ gives the lower endpoint in
\eqref{graph:first-interval}. The interval is nonempty because
$V\le A+C$, $V\le A+\alpha D$, and all capacities are nonnegative.

Conversely, for any $z$ in this interval, take
\[
 f_{A_LB_R}=(V-\alpha D)_+,\qquad
 f_{A_LD_R}=V-z-(V-\alpha D)_+,\qquad f_{C_LD_R}=z.
\]
These quantities are nonnegative: if $V>\alpha D$, use $z\le\alpha D$;
otherwise use $z\le V$. The $B_R$ use is at most $\alpha B$ since
$V\le\alpha(B+D)$, and the $D_R$ use is
$V-(V-\alpha D)_+\le\alpha D$. Their left uses are $V-z\le A$ and
$z\le C$. This realizes every point of the interval and proves that $V$
is the maximum. Reflecting the path, that is, exchanging $A$ with $D$ and
$B$ with $C$, proves the second assertion. Explicitly, with $z=w_B$ its
flows are $(V'-\alpha A)_+$ on $D_LC_R$, $V'-z-(V'-\alpha A)_+$ on
$D_LA_R$, and $z$ on $B_LA_R$.
\end{proof}

Define $\operatorname{clamp}(x;\ell,h)=\min\{h,\max\{\ell,x\}\}$ and
select
\begin{align}
 w_C&=\operatorname{clamp}\bigl(\beta C;
             \max\{0,V-A\},\min\{C,\alpha D,V\}\bigr),
 &w_A&=V-w_C,\label{graph:selector}\\
 w_B&=\operatorname{clamp}\bigl(\beta B;
             \max\{0,V'-D\},\min\{B,\alpha A,V'\}\bigr),
 &w_D&=V'-w_B.\nonumber
\end{align}
The vector $w$ is continuous, nonnegative, homogeneous, and piecewise
linear, with $0\le w_v\le c_v$. The following equivalent formulas expose
its linear cells:
\begin{align}
 w_A&=\min\{A,\alpha B+(\alpha D-\beta C)_+\},
 &w_D&=\min\{D,\alpha C+(\alpha A-\beta B)_+\},\label{graph:explicit-uses}\\
 w_C&=\min\{C,\alpha D,\max\{\beta C,\alpha(B+D)-A\}\},
 &w_B&=\min\{B,\alpha A,\max\{\beta B,\alpha(A+C)-D\}\}.\nonumber
\end{align}
To verify the first orientation algebraically, put
$U=\min\{C,\alpha D\}$, $T=\alpha(B+D)$, $t=\beta C$, and
$m=\min\{U,t\}=\min\{\alpha D,\beta C\}$.
Then $T\ge U$ and $V=\min\{T,A+U\}\ge U$, so the upper endpoint is
$U$. Since $V-A\le U$, the clamp equals
\[
 w_C=\min\{U,\max\{t,V-A\}\}
     =\max\{m,V-A\}
     =\min\{U,\max\{t,T-A\}\}.
\]
Consequently
$w_A=\min\{A,V-m\}=\min\{A,T-m\}$, giving
\eqref{graph:explicit-uses}. Reflection gives the other two formulas.

Set
\begin{equation}\label{graph:prices}
 p_v(c)=\begin{cases}w_v(c)/c_v,&c_v>0,\\0,&c_v=0,\end{cases}
 \qquad y_i(b)=\sum_{v\in\{A,B,C,D\}}b_{iv}p_v\left(\sum_j b_j\right).
\end{equation}
Prices lie in $[0,1]$. The price chosen at an absent coordinate has no
effect on any payment.

The selected left uses are realized by the flows in
Lemma~\ref{graph:flow-interval}. Collapse both directed copies of each
path edge into one edge flow. At vertex $v$, its incident use is at most
$c_v+\alpha c_v$, from the left and right copies respectively. Dividing
by $1+\alpha$ yields a feasible original packing. Therefore
\begin{equation}\label{graph:flow-funding}
 \sum_i y_i(b)=\sum_v w_v(c)=V+V'\le(1+\alpha)F_G(c).
\end{equation}
This argument uses fractional flows and no integrality assumption.

\subsection{The restoration inequalities on every smooth cell}

At a strictly positive aggregate $c$, let $a$ be a coalition's current
aggregate holding, with $0\le a\le c$, and write
$\theta_v=a_v/c_v\in[0,1]$. On a smooth cell define
\begin{equation}\label{graph:restoration-coefficients}
 g_v=p_v+\sum_j a_j\partial_v p_j.
\end{equation}
This differentiates the coalition's total payment when its holding and
the aggregate capacity are restored together, while outsiders stay fixed.
We prove
\begin{equation}\label{graph:cover-certificate}
 g_v\ge0\quad\text{for every vertex},\qquad
 g_A+g_B\ge1,\quad g_A+g_D\ge1,\quad g_D+g_C\ge1.
\end{equation}
Ties and zero capacities will be handled by direct integration below.

First record the single-edge calculation used in the first cells. For
edge capacities $X\ge Y>0$, put $t=Y/X$ and take
$w_X=\alpha Y$, $w_Y=\min\{Y,\alpha X\}$.
When $t\le\alpha$, the restoration coefficients are
\[
 g_X=\alpha t(1-\theta_X),\qquad g_Y=1+\alpha\theta_X.
\]
They are nonnegative and their sum is
$1+\alpha t+\alpha(1-t)\theta_X\ge1$.
When $t\ge\alpha$, they are
\[
 g_X=\alpha t(1-\theta_X)+\alpha\theta_Y,\qquad
 g_Y=(\alpha/t)(1-\theta_Y)+\alpha\theta_X.
\]
They are nonnegative. Minimizing their sum over the holding fractions gives
$\alpha(1+t)\ge\alpha(1+\alpha)>1$: the coefficient of $\theta_X$ in
the sum is nonnegative, and that of $\theta_Y$ is nonpositive.
Swapping $X,Y$ handles the other capacity ordering.

For the path put
\[
 u=A-\alpha B,\qquad v=D-\alpha C.
\]
The following cases exhaust the smooth cells; the displayed weak
inequalities describe their closures.

\paragraph{A nonpositive difference.}
If $u\le0$, formula \eqref{graph:explicit-uses} reduces to
\[
 w_A=A,\quad w_B=\alpha A,\quad
 w_C=\min\{C,\alpha D\},\quad w_D=\min\{D,\alpha C\}.
\]
For the first pair,
\[
 g_A=1+\alpha\theta_B\ge1,\qquad
 g_B=(\alpha A/B)(1-\theta_B)\ge0.
\]
The other pair satisfies the single-edge calculation, so $g_C,g_D\ge0$
and $g_C+g_D\ge1$. The cross edge is covered by $g_A\ge1$.
If $v\le0$, reflection gives the same conclusions with $g_D\ge1$.

For the remaining cells $u,v>0$, and the common formulas are
\begin{align*}
 w_A&=\alpha B+\min\{u,\alpha v\},&
 w_D&=\alpha C+\min\{v,\alpha u\},\\
 w_B&=\min\{B,\beta B+(\alpha u-v)_+\},&
 w_C&=\min\{C,\beta C+(\alpha v-u)_+\}.
\end{align*}
Since $\alpha<1$, at most one of $\alpha u-v$ and $\alpha v-u$ is
positive.

\paragraph{The middle cone: $\alpha v\le u$ and $\alpha u\le v$.}
Here
\[
 w_B=\beta B,\quad w_C=\beta C,\quad
 w_A=\alpha(B+D)-\beta C,\quad w_D=\alpha(A+C)-\beta B.
\]
Moreover $p_A,p_D\ge\beta$. For example,
$\alpha u\le v$ gives
$w_A=\alpha B+\alpha v\ge\beta(\alpha B+u)=\beta A$;
the other inequality follows by reflection. Direct differentiation yields
\begin{align*}
 g_A&=p_A(1-\theta_A)+\alpha\theta_D,&
 g_D&=p_D(1-\theta_D)+\alpha\theta_A,\\
 g_B&=\beta+\alpha\theta_A-\beta\theta_D,&
 g_C&=\beta+\alpha\theta_D-\beta\theta_A.
\end{align*}
All are nonnegative. Their edge-sum minima over $\theta\in[0,1]^4$ are
\[
 \min(g_A+g_B)=\beta+\min\{p_A,\alpha\},\qquad
 \min(g_D+g_C)=\beta+\min\{p_D,\alpha\},
\]
and
\[
 \min(g_A+g_D)=\min\{p_A,\alpha\}+\min\{p_D,\alpha\}.
\]
Every displayed value is at least $2\beta=1$.

\paragraph{An outer cone: $u\le\alpha v$.}
The opposite outer cone is its reflection. In this cone
\[
 w_A=A,\quad w_B=\beta B,\quad
 w_D=\alpha(A+C)-\beta B,\quad
 w_C=\min\{C,\alpha(B+D)-A\}.
\]
If $w_C<C$, then $p_C\ge\beta$ because $\alpha v-u\ge0$.
The unsaturated condition gives
$\alpha v-u\le(1-\beta)C$. Multiplying by $\alpha$ and using
$\beta=\alpha^2$ gives
\[
 \beta v\le\alpha u+\alpha(1-\beta)C,
 \qquad \beta D\le\alpha u+\alpha C=w_D.
\]
Thus $p_D\ge\beta$ as well. The coefficients are
\begin{align*}
 g_A&=1-\theta_C+\alpha\theta_D,&
 g_B&=\beta+\alpha\theta_C-\beta\theta_D,\\
 g_C&=p_C(1-\theta_C)+\alpha\theta_D,&
 g_D&=p_D(1-\theta_D)+\alpha\theta_C.
\end{align*}
They are nonnegative. Direct affine minimization gives
\begin{align*}
 \min(g_A+g_B)&=\alpha+\beta\ge1,\\
 \min(g_C+g_D)&=\min\{p_C,\alpha\}+\min\{p_D,\alpha\}\ge2\beta=1,\\
 \min(g_A+g_D)&=\alpha+\min\{p_D,\alpha\}\ge\alpha+\beta\ge1.
\end{align*}
For example the first sum is
$1+\beta+(\alpha-1)\theta_C+(\alpha-\beta)\theta_D$, whose minimum
is $\alpha+\beta$.

If instead $w_C=C$ on a smooth saturated cell, the coefficients are
\[
 g_A=1+\alpha\theta_D,\qquad g_C=1+\alpha\theta_D,
 \qquad g_B=\beta(1-\theta_D),\qquad
 g_D=p_D(1-\theta_D).
\]
They are nonnegative; $g_A\ge1$ covers $AB$ and $AD$, and $g_C\ge1$
covers $CD$. This completes \eqref{graph:cover-certificate} on every
smooth cell for every holding fraction, without restricting mixed ownership.

\subsection{Finite restorations, ties, and zero reports}

For outsider aggregate $o\ge0$ and coalition aggregate $a\ge0$, define
\[
 H(o,a)=\sum_v a_v p_v(o+a).
\]
This function is jointly continuous on its entire nonnegative domain.
At a coordinate whose aggregate tends to zero, its contribution lies
between zero and $a_v$ and hence tends to zero; at a positive aggregate
coordinate, continuity follows from that of $w$. Also $H(o,0)=0$.

Fix $0\le a'\le a$ and put $d=a-a'$. The finitely many polyhedral cells
of $w$ have boundaries contained in finitely many proper hyperplanes.
Choose a strictly positive vector $z$ having nonzero inner product with
each of their nonzero normals. Such a vector exists because a finite union
of proper hyperplanes cannot cover the positive orthant. Choose
$\varepsilon\downarrow0$ omitting the finitely many positive values for
which the segment
\[
 c(t)=o+\varepsilon z+a'+td,\qquad 0\le t\le1,
\]
is contained in any boundary hyperplane. Indeed, containment requires the
hyperplane to be parallel to $d$ and imposes one linear equation in
$\varepsilon$ with a nonzero coefficient. All other boundary intersections
occur at at most one parameter value per hyperplane.

The segment is strictly positive. Its coalition payment
$H(o+\varepsilon z,a'+td)$ is continuous and piecewise continuously
differentiable, with only finitely many breakpoints. Away from them its
derivative is $g\cdot d$, and its holding fractions are
$(a'_v+td_v)/c_v(t)\in[0,1]$. The verified inequalities therefore apply.
For any edge packing $x$ feasible at $d$, they imply
\[
 \sum_e x_e\le\sum_{e=uv}(g_u+g_v)x_e
 =\sum_v g_v\sum_{e\ni v}x_e\le\sum_v g_vd_v.
\]
Taking the maximum gives $g\cdot d\ge F_G(d)$. Integrating on the
finitely many pieces, using continuity at their endpoints, yields
\[
 H(o+\varepsilon z,a)-H(o+\varepsilon z,a')\ge F_G(d).
\]
Joint continuity as $\varepsilon\downarrow0$ proves
\begin{equation}\label{graph:finite-out}
 H(o,a)-H(o,a')\ge F_G(a-a')
 \qquad(o\ge0,\ 0\le a'\le a).
\end{equation}
The perturbation is an analytic approximation to the globally defined
formula $H$; no perturbed owner profile is required to fit the fixed boxes,
and no owner is added or reassigned in the asserted mechanism.

For an actual coalition $T$, take $o=\sum_{i\notin T}b_i$,
$a=\sum_{i\in T}b_i$, and $a'=\sum_{i\in T}b'_i$ for its binding
underreports. Its aggregate payment is exactly $H(o,a)$, so
\eqref{graph:finite-out} is full OUT for the original owners, including
passive members. It includes segments on ties, zero capacities, the
all-zero profile, and the grand coalition. Taking $a'=0$ proves coverage.
Together with \eqref{graph:flow-funding}, this proves
Theorem~\ref{graph:theorem} for the path.

\subsection{Complete blow-ups and disjoint unions}

Let $X\in\{A,B,C,D\}$ now denote one of four independent vertex groups,
and aggregate capacities as $c_X=\sum_{v\in X}c_v$. The production value
of a complete blow-up equals the path value at these four aggregates.
Aggregating any original packing proves one inequality. To prove the
other, let $f_{XY}$ be a feasible aggregate path packing. For $c_X>0$ assign
the incidence of this group-pair flow to $v\in X$ by
\[
 r_{v,XY}=\frac{c_v f_{XY}}{c_X}.
\]
The sum over incident group pairs is at most $c_v$, and
$\sum_{v\in X}r_{v,XY}=f_{XY}$. If $f_{XY}>0$, both group capacities
are positive and we may put
\[
 x_{vw}=\frac{r_{v,XY}r_{w,XY}}{f_{XY}}
 \qquad(v\in X,\ w\in Y)
\]
on every edge of that complete bipartite pair. Its two marginals are the
assigned incidences, and its total is $f_{XY}$. Zero-flow pairs receive
zero flow. Zero-capacity groups have no positive incident aggregate flow,
so no division by zero is used. This constructs an original feasible
packing of exactly the aggregate objective.

Apply the four path prices to every resource in its corresponding group.
Each actual owner's payment equals the path-rule payment on its own
aggregated report. Every actual withholding maps to a coordinatewise
withholding of that same owner's aggregate, and the equality of production
values holds for all coalition and withheld bundles. Thus funding,
coverage, and full OUT transfer exactly, retaining the original mixed
owners and their fixed endowment boxes.

For a finite disjoint union, sum the component payments. Both production
and withheld production are additive across components, so summing their
OUT, coverage, and funding inequalities preserves the constant $1+\alpha$.
An owner may hold resources in several components. Isolated resources
contribute no production and may receive zero payment. This completes
Theorem~\ref{graph:theorem}.

\begin{remark}[Scope of the funding number]\label{graph:scope}
At $(A,B,C,D)=(4,1,1,1)$, the path value is two and
$(V,V')=(2\alpha,2)$, so this particular rule attains the ratio
$1+\alpha$. This is not a lower bound for arbitrary mechanisms on the
path. The proof does not extend the guarantee to arbitrary edge profits
or to incomplete connections between the four groups, and makes no
claim of historical novelty.
\end{remark}


\subsection{An obstruction to a direct leaf-layer recursion}
The path theorem does not license independent leaf reservations throughout a
tree. The following obstruction concerns a specified selection procedure;
it is not a lower bound for arbitrary mechanisms on trees.

Fix $0<\alpha<1$ and $0\le\beta<1$. In each orientation of the scaled
bipartite double cover, process left vertices in a leaf-layer order of the
original tree: all original leaves precede newly exposed leaves, and each
successive layer is exhausted before the next. Among maximum flows consistent
with previous choices, clamp each current left use to $\beta c_v$ within its
feasible interval. Pay owners by their resource fractions of these left uses,
as in \eqref{graph:prices}. On P4, $\alpha=1/\sqrt2$ and $\beta=1/2$
recover \eqref{graph:selector}. The collapsed-flow funding argument still
applies on every tree, but incentives need not.

\begin{proposition}[Failure of fixed leaf reserves]\label{graph:recursion-failure}
For every $0<\alpha<1$ and $0\le\beta<1$, this leaf-layer selector fails
full OUT on P6, already for a singleton coalition consisting of one mixed owner.
The failure is independent of tie-breaking within the first leaf layer and
of every selection choice after the first leaf in each orientation is fixed.
\end{proposition}
\begin{proof}
Use the path $0-1-2-3-4-5$, choose
\[
 0<q<\min\{1-\beta,\alpha^{-2}-1,1\},
 \qquad c(q)=(1,\alpha^{-1},q,q,\alpha^{-1},1).
\]
One fixed owner holds all resources at $0,2,3,5$; two outsiders hold resources
$1,4$ respectively. The first owner alone reduces both middle reports from
$q$ to $q/2$, leaving its two endpoint reports unchanged. Its withheld bundle
has production $q/2$ on edge $2-3$.

Consider the orientation with left vertices $\{0,2,4\}$. The flows
\[
 f_{0_L1_R}=1-q,\quad f_{2_L1_R}=q,\quad
 f_{4_L3_R}=\alpha q,\quad f_{4_L5_R}=\alpha
\]
are feasible: in particular $\alpha(1+q)<\alpha^{-1}$.
They saturate every right capacity, so their value $1+\alpha(1+q)$ is
maximum. Every maximum flow saturates $1_R$, whose only incoming sources
are $0_L,2_L$. Since the latter has capacity $q$, its original leaf satisfies
$w_0\ge1-q$. The displayed flow attains this lower endpoint. As
$1-q>\beta c_0=\beta$, the first leaf clamp fixes $w_0=1-q$.
Saturating $1_R$ then forces source $2_L$ to spend its entire $q$ there;
all other right use must come from $4_L$. Thus every remaining left use
is forced. Reflection proves the same statement for the other orientation,
whose first original leaf is $5_L$. Hence
\[
 w(q)=(1-q,\alpha(1+q),q,q,\alpha(1+q),1-q).
\]
The same argument applies at $q/2$, since all strict inequalities persist.
The deviating owner's total payment is
\[
 (1-q)+q+q+(1-q)=2
\]
at both reports. Its payment loss is zero, contradicting the required
outside-production compensation $q/2>0$.
\end{proof}

For the theorem's constants, $q=1/4$ gives the concrete capacities
$(1,\sqrt2,1/4,1/4,\sqrt2,1)$ and a productive withdrawal of $1/8$
at each middle vertex. Retuning a fixed fraction below one or taking
$\alpha$ closer to one cannot repair this leaf-layer architecture. The
proposition does not cover reserve fraction one, report-dependent reserves,
other global selectors, or orders that process a newly exposed vertex before
an original leaf in its orientation. In particular it establishes no
impossibility for arbitrary full-OUT tree mechanisms. A viable recursion
must retain additional information about the simultaneous payment response
across the central edge.

\section{Interpretation, scope and remaining questions}\label{sec:discussion}
The two parameters answer different questions. The dimension frontier allows
arbitrary technologies on $d$ resources and protects every coalition. The
sparse frontier limits each activity to $r$ resources, allows $d$ to grow,
and protects coalitions of at most $s$ owners while retaining coverage for all
coalitions. A resource-dimension upper cannot be substituted for an
owner-independent bound depending only on $r,s$.

At arity three, the square-root-to-quadratic gap remains open. The lower uses
exponentially many resource rows. Pure ownership is possible, but requires
more owners than the fully eligible construction; Appendix~\ref{sec:arity}
keeps the counts separate. Neither the exponent nor its sharper constant
claims minimal witness size. The simple binomial upper evaluates finitely
many coalition production programs, potentially exponentially many. Efficient
approximate evaluation requires separate incentive and coverage analysis.

The graph construction in Section~\ref{graph:section} protects every coalition
while allowing the number of resources to grow within complete P4 blow-ups and
disjoint unions. Its factor is uniform in those counts, but its interaction
structure is restricted. Extending the selector to arbitrary trees or graphs,
and determining the optimal factor even on P4, remain open within this program.
The extra funding is $F/\sqrt2$, not $(1+1/\sqrt2)F$ in addition to output.

The proposed contribution is the quantitative link between fixed-owner
coalition withholding and funding, particularly the projective cap exchange
that supplies a constant seed without a dummy-zero assumption. Projective
geometry, positive hypergraph supports, semivalues, and coalition monotonicity
are established ingredients. A specialist must assess whether a prior theorem,
possibly under different terminology or normalization, already implies the
present bounds. The accompanying priority matrix records inspected versions
and unresolved original-source comparisons. Internal model reviews and exact
certificates are not a historical-novelty certificate or human peer review.

A conditional infrastructure interpretation would give each owner genuine
rights over compute inputs and firm power capacity, with physically feasible
partial withdrawals and a sponsor who benefits from the stronger contract.
Public telemetry does not establish those rights. The earlier bounded
feasibility check supplied no compatible held-out outcome panel or sponsor
surplus estimate. Moreover, a Gaussian reserve model is not automatically an
arity-three packing technology. We therefore make no measured pooling,
reliability, energy-saving or net-value claim here.

This manuscript concerns coordinated withholding. It does not settle the
separate owner-count/resource-dimension gap for individual incentives in the
companion truthful-production project. Optional floor parameters, structured
production refinements beyond the explicitly proved graph class, and application
measurement studies are not separate new claims of this draft. The immediate questions are correctness of the
fixed-institution reductions, priority of the quantitative bounds, and whether
these results form a useful unit of publication.

\appendix
\section{Fixed arity and pure ownership}
\label{sec:arity}
\label{arity:appendix}

The projective construction of Theorem~\ref{thm:projective} extends to every
fixed arity $r\geq 3$ by
replacing binary triples with binary circuits of size $r$.  This appendix
gives the additional geometry and records the exact finite bounds.  It also
gives a separate realization with one pure owner per resource.  The resource
count is exponential in the length of the construction in both realizations;
the pure realization has exponentially many owners as well.

Throughout, payments are arbitrary nonnegative functions on the entire fixed
real report domain.  For a report profile $b=(b_i)_{i\in N}$, write
$b_C=\sum_{i\in C}b_i$.  The hypotheses used below are coalition coverage and
a uniform funding bound,
\begin{equation}
  \sum_{i\in C}y_i(b)\geq F(b_C)\quad(C\subseteq N),
  \qquad
  \sum_{i\in N}y_i(b)\leq\kappa F(b_N),
  \label{arity:contract}
\end{equation}
at every feasible profile, together with size-$s$ coalition withholding
monotonicity: if $|C|\leq s$, $b'_i\geq b_i$ coordinatewise for $i\in C$,
and $b'_i=b_i$ for $i\notin C$, then
\begin{equation}
  \sum_{i\in C}y_i(b')\geq\sum_{i\in C}y_i(b).
  \label{arity:tcwm}
\end{equation}
Members of $C$ may have unchanged reports.  Coverage in
\eqref{arity:contract} is required for \emph{every} coalition, including
coalitions larger than $s$.  No continuity, formula for payments, anonymity,
or zero payment to a zero reporter is assumed.  Since full OUT implies
\eqref{arity:tcwm}, every lower bound below also applies under full OUT.

\subsection{The binary-circuit construction}
\label{arity:construction-section}

Fix integers $k\geq2$ and $M\geq1$, and put $r=k+1$.  Let
\[
  V=\mathbb F_2^{\,1+(k-1)M}
\]
have the specified basis
\[
  v,\qquad e_{j,\ell}\quad
  (1\leq j\leq M,\ 1\leq\ell\leq k-1).
\]
Resource rows are indexed by $V\setminus\{0\}$, so their number is
\begin{equation}
  d=2^{1+(k-1)M}-1.
  \label{arity:row-count}
\end{equation}
A binary circuit means a nonempty set of distinct rows whose sum over
$\mathbb F_2$ is zero and none of whose proper nonempty subsets has sum zero.
The activities are all binary circuits of cardinality $k+1$.  Each activity
has profit one and consumes one unit of each of its $k+1$ rows.  Thus, with
$\mathcal A$ denoting the set of activities, the full technology is
\begin{equation}
  F(c)=\max\left\{
    \sum_{A\in\mathcal A}x_A:
    x_A\geq0,\quad
    \sum_{\substack{A\in\mathcal A\\u\in A}}x_A\leq c_u
       \text{ for every }u\in V\setminus\{0\}
  \right\}.
  \label{arity:technology}
\end{equation}
Every activity has arity exactly $r$, and every $g\in\operatorname{GL}(V)$
permutes the resources and the activities of this finite technology.

For the first realization take the fixed actual owner set
\[
  N=\{p\}\ \cup\
  \{h_{j,\ell}:1\leq j\leq M,\ 1\leq\ell\leq k-1\}
  \ \cup\ \{\ell_j:1\leq j\leq M\}.
\]
There are $n=1+kM$ owners.  Every owner has the all-ones endowment and the
same fixed report cube $[0,1]^d$.  Let $\delta_u\in\R^d$ denote one unit of
resource $u$, to distinguish real resource quantities from the binary
vectors that index resources.  Set
\[
  w_j=v+\sum_{\ell=1}^{k-1}e_{j,\ell}\quad\text{in }V.
\]
Basis independence makes all displayed root, high and low rows nonzero and
pairwise distinct.
For $1\leq m\leq M$, define a feasible profile $C_m$ by
\begin{equation}
  b_p=\delta_v,\qquad
  b_{h_{j,\ell}}=\delta_{e_{j,\ell}},\qquad
  b_{\ell_j}=
  \begin{cases}
    \delta_{w_j},&j\leq m,\\
    0,&j>m.
  \end{cases}
  \label{arity:profiles}
\end{equation}
All actual owners and all their report domains remain fixed as $m$ changes.
In particular, the high owners of inactive branches already report their
basis rows.

\begin{lemma}[Supported activities]
\label{arity:supported}
At $C_m$, the only activities whose rows all have positive capacity are
\[
  A_j=\{v,e_{j,1},\ldots,e_{j,k-1},w_j\},
  \qquad 1\leq j\leq m.
\]
Consequently $F(b_N(C_m))=1$.
\end{lemma}

\begin{proof}
Consider a nonempty zero-sum subset of the larger set
\[
  S=\{v\}\cup\{e_{j,\ell}:j,\ell\}\cup\{w_j:1\leq j\leq M\}.
\]
Let $J$ be its set of selected low-row indices and put $t=|J|$.  The basis
coordinate $e_{j,\ell}$ occurs in only two members of $S$: the row
$e_{j,\ell}$ itself and the row $w_j$.  Cancellation therefore requires all
$k-1$ high rows of every branch in $J$, and no high rows of any other branch.
Cancellation of the $v$ coordinate requires the root row precisely when $t$
is odd.  If $t=0$, these conditions give the empty subset.  Every nonempty
zero-sum subset thus has cardinality
\begin{equation}
  kt+(t\bmod 2).
  \label{arity:dependency-cardinality}
\end{equation}
For $k\geq2$, this cardinality equals $k+1$ precisely when $t=1$: if $t\geq2$
it is at least $2k>k+1$.  The corresponding set is $A_j$, and it is minimal
because its $k$ rows $v,e_{j,1},\ldots,e_{j,k-1}$ are independent and $w_j$
is their sum.  At $C_m$ its low row has positive capacity exactly when
$j\leq m$, which proves the classification.

Running any one active branch at level one is feasible.  Every supported
activity consumes the root row, whose capacity is one, so production cannot
exceed one.  Equivalently, a dual certificate for the \emph{full} technology
assigns price one to $v$ and to every zero-capacity row, and price zero to
the other rows.  A supported activity contains $v$; every other activity
contains a zero-capacity row.  Each dual constraint is therefore satisfied,
and the dual objective is one.
\end{proof}

\begin{lemma}[Seed and private-owner orbits]
\label{arity:orbits}
After lawful finite averaging over resource automorphisms and owner
permutations, the owners of the first branch at $C_1$ have a common payment
$q_1\geq1/(k+1)$.  At $C_m$, all $km$ active private owners
\[
  \{h_{j,\ell}:j\leq m,\ 1\leq\ell\leq k-1\}
  \cup\{\ell_j:j\leq m\}
\]
have a common payment $q_m$.
\end{lemma}

\begin{proof}
The finite averaging argument of Lemma~\ref{lem:average} applies to
$\operatorname{GL}(V)\times\mathfrak S_N$.  A binary linear automorphism
acts on a real report by \emph{permuting its resource coordinates}; it does
not take binary linear combinations of real capacities.  These resource
permutations preserve \eqref{arity:technology}, the real report cubes and
coordinatewise order.  Owner permutations preserve the common fixed report
domains, coalition sizes and the coverage and comparison inequalities.
Consequently conjugating a rule and averaging preserves nonnegativity,
\eqref{arity:contract}, \eqref{arity:tcwm}, and its funding factor.  It
suffices to prove the lower bound for this equivariant averaged rule.

At $C_1$, any permutation of the $k$ basis vectors
$v,e_{1,1},\ldots,e_{1,k-1}$ fixes their sum $w_1$ and fixes every high
basis row of an inactive branch.  Also the map
\[
  v\longmapsto w_1,\qquad e_{j,\ell}\longmapsto e_{j,\ell}
       \quad\text{for all }j,\ell
\]
is invertible: the coefficient of $v$ in the replacement for $v$ is one.
It exchanges $v$ and $w_1$.  These maps preserve the entire multiset of
reported rows at $C_1$, including the inactive high reports and zero low
reports.  Combined with the matching owner permutations, they stabilize
$C_1$ and act transitively on its $k+1$ first-branch owners.  Their coalition
can run $A_1$ at level one, so coverage gives $(k+1)q_1\geq1$.  This use of
coverage does not require $k+1\leq s$.

At $C_m$, permute the active branch blocks, or permute the high basis vectors
within one active branch, fixing $v$ and the inactive basis vectors.
Additionally, for each $j\leq m$ the basis replacement
\[
  e_{j,1}\longmapsto w_j
\]
with every other basis vector fixed is invertible, since the coefficient of
$e_{j,1}$ in $w_j$ is one.  It fixes $v$, exchanges $e_{j,1}$ with $w_j$,
and fixes all other displayed rows except these two.  Within-branch
permutations and this exchange connect all $k$ private rows of a branch;
active branch permutations connect the branches.  They preserve the entire
profile, so the induced stabilizers are transitive on all $km$ active
private owners.  Equivariance gives their common payment $q_m$.
\end{proof}

\begin{theorem}[Finite fixed-arity lower bound]
\label{arity:finite-bound}
Fix integers $k\geq2$, $s\geq2$ and $M\geq1$.  There is a finite institution
of arity exactly $k+1$, with $1+kM$ fixed fully eligible owners and
$2^{1+(k-1)M}-1$ resource rows, for which every payment rule satisfying
\eqref{arity:contract}--\eqref{arity:tcwm} has funding factor at least
\begin{equation}
  L_{k,M,s}
  :=1+\frac{k(M-1)}{k+1}P_{k,M,s},\qquad
  P_{k,M,s}:=\prod_{m=2}^{M}\frac{a_m}{a_m+1},\qquad
  a_m:=\min\{s-1,k(m-1)\}.
  \label{arity:finite-tail}
\end{equation}
Empty products are one.  If $s\geq k(M-1)+1$, including $s=1,M=1$, the
same construction gives
\begin{equation}
  \kappa\geq\frac{k}{k+1}Q_{k,M},\qquad
  \kappa\geq1+\frac{k(M-1)}{(k+1)M}Q_{k,M},\qquad
  Q_{k,M}:=\prod_{m=2}^{M}\frac{km}{k(m-1)+1}.
  \label{arity:untruncated}
\end{equation}
\end{theorem}

\begin{proof}
Use \eqref{arity:profiles} and the averaged rule from
Lemma~\ref{arity:orbits}.  In the transition from $C_{m-1}$ to $C_m$, only
the same actual owner $\ell_m$ changes its report, from $0$ to $\delta_{w_m}$.
Select any $a_m$ of the $k(m-1)$ old active private owners and protect them
together with $\ell_m$.  This fixed coalition has $a_m+1\leq s$ members.
All selected old members keep their reports unchanged.  At the old profile
their total payment is $a_mq_{m-1}$, and the restoring owner's old payment
is nonnegative.  At the new profile all selected members, including the
restorer, receive $q_m$.  Hence \eqref{arity:tcwm} implies the exact recurrence
\begin{equation}
  (a_m+1)q_m\geq a_mq_{m-1},\qquad 2\leq m\leq M.
  \label{arity:recurrence}
\end{equation}
The selected old owners need not comprise whole branches.  Iterating
\eqref{arity:recurrence} from $(k+1)q_1\geq1$ gives
\begin{equation}
  q_M\geq\frac{P_{k,M,s}}{k+1}.
  \label{arity:terminal-q}
\end{equation}

Write $z$ for the root owner's payment at $C_M$.  Coverage for the actual
owners of any complete branch and the budget at $C_M$ give respectively
\[
  z+kq_M\geq1,\qquad
  \kappa\geq z+kMq_M.
\]
The second inequality uses Lemma~\ref{arity:supported} and discards every
other owner's payment only by nonnegativity.  Combining them yields
$\kappa\geq1+k(M-1)q_M$, proving \eqref{arity:finite-tail}.  This includes
$M=1$, when terminal coverage alone gives $\kappa\geq1$.

If $s\geq k(M-1)+1$, protect \emph{all} old active private owners at every
transition.  The recurrence becomes
\[
  [k(m-1)+1]q_m\geq k(m-1)q_{m-1}.
\]
Writing $T_m=kmq_m$ gives
\[
  T_1\geq\frac{k}{k+1},\qquad
  T_m\geq\frac{km}{k(m-1)+1}T_{m-1}.
\]
Thus $T_M\geq kQ_{k,M}/(k+1)$, and the budget proves the first bound in
\eqref{arity:untruncated}.  Since
$q_M\geq Q_{k,M}/[(k+1)M]$, terminal branch coverage proves the second.
When $M=1$ there are no comparisons, so this argument also applies at $s=1$.

All comparisons just used occur inside the one original fixed institution.
Symmetry is a proof reduction, not a restriction on admissible payment
rules.  If the original rule is asymmetric, an expensive profile may be an
image of $C_M$ under the finite symmetry group, rather than $C_M$ itself.
The same lower bound on its worst-case funding factor follows.
\end{proof}

The finite bound also has a direct nonnegative rational certificate.
Put $c=k(M-1)$ and $P=P_{k,M,s}$.  Give the terminal coverage inequality
$z+kq_M\geq1$ and the budget inequality
$\kappa-z-kMq_M\geq0$ weight one each.  Give the seed inequality
$(k+1)q_1\geq1$ weight $cP/(k+1)$, and give recurrence $m$ weight
\begin{equation}
  \lambda_m=\frac{c}{a_m+1}
    \prod_{j=m+1}^{M}\frac{a_j}{a_j+1}.
  \label{arity:certificate-weights}
\end{equation}
For $M\geq2$, the coefficient of $q_M$ from the final two inequalities is
$-c$, canceled by $\lambda_M(a_M+1)$.  For $2\leq m<M$,
$\lambda_m(a_m+1)=\lambda_{m+1}a_{m+1}$, and
$\lambda_2a_2=cP$ cancels the seed term.  The $z$ coefficients also cancel,
leaving exactly $\kappa\geq1+cP/(k+1)$.  At $M=1$, $c=0$ and terminal
coverage plus budget already give the conclusion.  These are consequences
for arbitrary payment variables; no marginal-payment representation is used.

\subsection{Length optimization and the fixed-arity exponent}
\label{arity:rate-section}

\begin{corollary}[Optimized length within the construction]
\label{arity:optimized}
For fixed integers $k\geq2$ and $s\geq2$, set
\[
  h=\left\lceil\frac{s}{k}\right\rceil,
  \qquad \rho=\frac{s-1}{s}.
\]
The largest value of $L_{k,M,s}$ over positive integer lengths $M$ is
attained exactly at $M=s$ and $M=s+1$, with common value
\begin{equation}
  B_{k,s}
  =1+\frac{k(s-1)}{(k+1)h}Q_{k,h}\rho^{\,s-h}.
  \label{arity:optimized-bound}
\end{equation}
The length $M=s$ realizes this bound with
$n=1+ks$ fully eligible owners and $d=2^{1+(k-1)s}-1$ rows.  For fixed $k$,
\begin{equation}
  B_{k,s}-1\sim
  \frac{k}{k+1}\Gamma\!\left(1+\frac1k\right)
       k^{1/k}\exp\!\left(-\frac{k-1}{k}\right)
       s^{(k-1)/k}
       \qquad(s\longrightarrow\infty).
  \label{arity:asymptotic}
\end{equation}
This optimizes the displayed certificate over its length; it does not
identify the optimal funding factor among all mechanisms.
\end{corollary}

\begin{proof}
The cutoff satisfies $k(h-1)\leq s-1$ and $kh\geq s$.  Through $m=h$ the
recurrence is untruncated, and for $m\geq h+1$ its factor is $\rho$.
Telescoping the factors $(m-1)/m$ therefore gives
\[
  P_{k,h,s}=\frac{Q_{k,h}}{h},\qquad
  P_{k,M,s}=\frac{Q_{k,h}}{h}\rho^{M-h}\quad(M\geq h).
\]
Let $I_M=(M-1)P_{k,M,s}$.  We have $I_1=0<I_2$.  If $1<M<h$,
\[
  \frac{I_{M+1}}{I_M}
  =\frac{kM^2}{(M-1)(kM+1)}>1,
\]
because the numerator minus the denominator is $(k-1)M+1>0$.  If
$M\geq h$ and $M>1$, the ratio is
\[
  \frac{I_{M+1}}{I_M}=\frac{M\rho}{M-1},
\]
which exceeds one for $M<s$, equals one for $M=s$, and is less than one
for $M>s$.  As $h\leq s$, this proves the two asserted maximizers, including
the boundary case $h=1$.  Substituting $M=s$ proves
\eqref{arity:optimized-bound}.

The gamma functional equation gives the exact identity
\[
  Q_{k,h}
  =\frac{\Gamma(h+1)\Gamma(1+1/k)}{\Gamma(h+1/k)}.
\]
Stirling's formula yields
$Q_{k,h}\sim\Gamma(1+1/k)h^{(k-1)/k}$.  Moreover $h\sim s/k$ and
$\rho^{s-h}\longrightarrow\exp(-(k-1)/k)$.
Inserting these limits into \eqref{arity:optimized-bound} proves
\eqref{arity:asymptotic}.
\end{proof}

\begin{corollary}[An all-$s$ bound]
\label{arity:all-s}
For every fixed integer $r\geq3$ and every integer $s\geq1$, writing
$k=r-1$, the unrestricted worst-case optimal funding factor satisfies
\begin{equation}
  K_r(s)\geq
  \frac{k}{k+1}\left(\frac{s}{2k}\right)^{(k-1)/k}.
  \label{arity:elementary-rate}
\end{equation}
In particular,
$K_r(s)=\Omega_r(s^{(r-2)/(r-1)})$.
\end{corollary}

\begin{proof}
Put $\beta=(k-1)/k$.  For $m\geq2$, concavity of $x^\beta$ gives
\[
  \frac{km}{k(m-1)+1}
    =\frac{m}{m-\beta}
    \geq1+\frac{\beta}{m}
    \geq\left(1+\frac1m\right)^\beta.
\]
Consequently $Q_{k,M}\geq((M+1)/2)^\beta$, also at $M=1$.
Choose $M=\lceil s/k\rceil$.  Then
$k(M-1)+1\leq s$ and $M\geq s/k$, so the first bound in
\eqref{arity:untruncated} implies \eqref{arity:elementary-rate}.
The finite institution is admissible for the worst-case class defining
$K_r(s)$.
\end{proof}

\subsection{A separate realization with pure owners}
\label{arity:pure-section}

\begin{corollary}[Fixed pure-owner realization]
\label{arity:pure}
All the lower bounds above also hold for institutions with one fixed pure
unit owner per resource row.  For the length-$M$ instance, this realization
has
\begin{equation}
  n=d=2^{1+(k-1)M}-1,
  \label{arity:pure-counts}
\end{equation}
and every owner's fixed report domain is the full real interval along its
own resource coordinate.  It does not have the fully eligible realization's
owner count $1+kM$.
\end{corollary}

\begin{proof}
Keep the technology \eqref{arity:technology}.  For each row
$u\in V\setminus\{0\}$ introduce the actual owner $i_u$ with fixed
endowment $\delta_u$ and report domain
\[
  \{t_u\delta_u:0\leq t_u\leq1\}.
\]
All these owners exist at every profile.  A resource automorphism $g$ acts
\emph{jointly} on resource rows and owner labels, sending owner $i_u$ to
owner $i_{g(u)}$.  On scalar reports the action is
\[
  (g\cdot t)_{g(u)}=t_u.
\]
This preserves the fixed institution: the transformed owner still reports
only its own resource.  The associated owner bijection takes a coalition
to another coalition of the same size, maps its total resource report by
the same technology automorphism, preserves coordinatewise restoration,
and leaves all outsiders unchanged in a comparison.  Conjugation therefore
preserves every coverage, funding and incentive inequality.  Averaging over
this finite joint $\operatorname{GL}(V)$ action is lawful.  No independent
owner permutation is required or used.

At the pure-owner version of $C_m$, set $t_u=1$ exactly for
\[
  S_m=\{v\}\cup\{e_{j,\ell}:j,\ell\}
                 \cup\{w_j:j\leq m\},
\]
and set $t_u=0$ for every other row.  Its aggregate capacities coincide with
those of \eqref{arity:profiles}, so Lemma~\ref{arity:supported}, including
the full-technology dual, applies without change.  Each seed or private
stabilizer described in Lemma~\ref{arity:orbits} preserves $S_m$ as a set.
Being a permutation of all nonzero rows, it also preserves its complement.
Its joint action thus fixes the entire report profile, including every
additional zero-report owner.  The same seed and private orbits give
$(k+1)q_1\geq1$ and a common payment $q_m$ for the $km$ active private owners.

In the transition from $C_{m-1}$ to $C_m$, only the same actual pure owner
$i_{w_m}$ restores its report from zero to one.  Protect this owner together
with the specified old active pure owners.  Their labels, eligibility and
reports remain fixed throughout the comparison.  Equation
\eqref{arity:recurrence} follows exactly as before, by discarding only the
restorer's nonnegative old payment.  Terminal branch coverage involves
the $k+1$ actual pure owners of that branch.  All omitted owners, including
additional zero reporters, have nonnegative payments, so the terminal
budget argument is unchanged.  This proves the finite bounds and all their
stated consequences with the counts \eqref{arity:pure-counts}.
\end{proof}

For example, $k=2$ specializes to the arity-three construction.  Its optimized
length-$s$ bound has coefficient
$\sqrt{2\pi}/(3\sqrt e)$ in front of $\sqrt s$, and $B_{2,6}=493/243$.
The length-six witness uses $127$ resource rows in either realization,
with $13$ fully eligible owners or $127$ pure owners.  Neither the
fixed-arity extension nor the change of ownership realization gives a
bound with a fixed resource dimension.

\begin{thebibliography}{99}
\bibitem[Bachrach et~al.(2009)]{BachrachEtAl2009}
Y. Bachrach, E. Elkind, R. Meir, D. Pasechnik, M. Zuckerman, J. Rothe,
and J. S. Rosenschein (2009).
The cost of stability in coalitional games.
\emph{SAGT 2009}, LNCS 5814, 122--134.
Inspected preprint: \url{https://arxiv.org/abs/0907.4385}.

\bibitem[Casajus(2011)]{Casajus2011}
A. Casajus (2011).
Collusion, symmetry, and the Banzhaf value.
Author manuscript, 31 August 2011.
\url{https://home.uni-leipzig.de/casajus/texts/colluba.pdf}.

\bibitem[Do\u{g}an(2013)]{Dogan2013}
E. Do\u{g}an (2013).
Absence-proofness: A new cooperative stability concept.
Working paper, 11 September 2013, inspected in the source audit.
Published as \emph{Absence-proofness: Group stability beyond the core}.
Published article: \emph{International Journal of Game Theory} 45, 601--616
(issue 2016; online publication 2015),
\url{https://doi.org/10.1007/s00182-015-0472-z}.
The published text has not been verified against the inspected version.

\bibitem[Ezra et~al.(2018)]{EzraEtAl2018}
T. Ezra, M. Feldman, E. Neyman, I. Talgam-Cohen and S. M. Weinberg (2018).
Complement-free couples must communicate: A hardness result for two-player
combinatorial auctions. \url{https://arxiv.org/abs/1811.09871}.

\bibitem[Feige et~al.(2014)]{FeigeEtAl2014}
U. Feige, M. Feldman, N. Immorlica, R. Izsak, B. Lucier and V. Syrgkanis (2014).
A unifying hierarchy of valuations with complements and substitutes.
\url{https://arxiv.org/abs/1408.1211}.

\bibitem[Immorlica et~al.(2008)]{ImmorlicaEtAl2008}
N. Immorlica, M. Mahdian and V. S. Mirrokni (2008).
Limitations of cross-monotonic cost-sharing schemes.
\emph{ACM Transactions on Algorithms} 4(2).
Inspected author full version: \url{https://immorlica.com/pubs/costshare.pdf}.

\bibitem[Meir et~al.(2011)]{MeirEtAl2011}
R. Meir, J. S. Rosenschein and E. Malizia (2011).
Subsidies, stability, and restricted cooperation in coalitional games.
\emph{IJCAI 2011}, 301--306.
\url{https://www.ijcai.org/Proceedings/11/Papers/060.pdf}.

\bibitem[Moulin and Sprumont(2004)]{MoulinSprumont2004}
H. Moulin and Y. Sprumont (2004).
On demand responsiveness in additive cost sharing.
Author version, 24 May 2004.
\url{https://www.ruf.rice.edu/~econ/papers/2004papers/dmdrespr41.pdf}.
Published in \emph{Journal of Economic Theory} 125(1), 1--35 (2005),
\url{https://doi.org/10.1016/j.jet.2004.05.003}.
Cited proposition numbering refers to the inspected author version.
\bibitem[Xiao and Fang(2021)]{XiaoFang2021}
H. Xiao and Q. Fang (2021).
Population monotonicity in matching games.
Inspected original: arXiv:2105.00621v1, especially Section~2 and Lemma~3.
\url{https://arxiv.org/abs/2105.00621}.

\end{thebibliography}

\end{document}
